% NAL 1644 — lecture modernisée du volume entier.
%
% Pass: Opus 5.5 (claude-opus-5-5), 2026-09-23 — first pass, unchecked against
% the views by a human, and an interpretation rather than a record. Where this
% file and the transcriptions differ in substance, the transcriptions
% (batch-01.fr.tex to batch-09.fr.tex) are the record.
%
% Sources read: the nine batch transcriptions of this volume, views 1–163, in
% full, twice — once for the mathematics, once for the apparatus (every
% \uncertain, \ill, \note and the absence of any \folio) — with their header
% comments and the coordinator's reconciliation block closing each header, and
% the NAL 1644 section of .claude/skills/transcribe/references/hand.md. The
% volume is completely transcribed: v. 4 and v. 146 are blank and not
% transcribed; every other view is. No image was consulted. No printed edition
% was consulted: not the Elzevier Opera of 1646, not any edition or study of
% Viète, of Tycho or of Copernicus. The later reader's page references to the
% 1646 edition (v. 1, 5, 69, 141) are reported as notes on the leaves and were
% not checked. Where the transcription has a gap, this reading says so.
%
% What the volume turned out to be (map settled by the reconciliation notes):
%
%   v. 1        cover, « Francisci Vietæ opera nonnulla », later references to
%               the 1646 edition and a table « 1 » / « 32 » (= ink leaves 1 and
%               32, where the two treatises begin: an observation, not a reading)
%   v. 2–4      small leaves A, B, C (C blank), another compact hand: sextic
%               equations reduced (« duplicata hypostasis », « climactica
%               paraplerosis »), a theorem on alternate terms
%   v. 5–68     « Francisci Vietæ De Recognitione Æquationum tractatus. »,
%               Cap. 1–XXI (heading on the leaf, in the text's hand)
%   v. 69–141   « Francisci Vietæ. De æquationum Emendatione tractatus
%               secundus. », Cap. I–XIIII, « Finis » (heading on the leaf)
%   v. 142–163  astronomy, no name: lunar prostaphaereses and their tables;
%               Tycho's lunar hypothesis « restored » (Tycho's name struck for
%               « ab Astroscopo Dano »); « Harmonici Ptolemaici Constructio
%               Caput XII », against Copernicus's geometry; outer planets and
%               Venus; Mercury. Breaks off mid-sentence.
%
%   One ink series of leaf numbers, 1–79 with « 5 bis », « 16 bis »,
%   « 46 bis », « 47 bis », crosses all three pieces; the writer cites it
%   (« fol. 18. », v. 43; « pag. 48 », v. 100). Not the library's foliation:
%   no folio is cited anywhere; sections cite views by \pagerange.
%   Chapter numbers are the leaf's corrected ones (VIII→IX, IX→X, XIIII→XV,
%   XVIII→XIX, XIX→XX in the first treatise; XX→X, XXI→XI … XXIIII→XIIII in
%   the second).
%
% Attribution. The only names for the authorship of the treatises are the two
% headings and the cover; the reading says what they say and identifies no
% hand. Doubtful notes that may mark the treatises as a copy are presented as
% doubtful: « transcripta ex exemplari » (v. 100, probable reading),
% « deerant hæc in exemplari » (v. 114, struck, very doubtful), and a third,
% « in exemplari excripto / vel potius adscriptum » (v. 20, partly read, a
% small pale hand), which batch 1 records and the brief did not list. The
% second hand of v. 13–18 speaks of its own work on angular sections
% (« demonstratum a nobis ad sectiones angulares »); reported, not identified.
% The astronomy is attributed to no one. \dating is empty, as in the catalogue;
% no date is inferred from the headings, from Tycho's name, or from the 1646
% references. Every « postérieur » is conditional on the headings'
% attribution, and is said so.
%
% Notation translated, each with a note at first use: A → x, E → y; the
% consonants kept, a magnitude named by its dimension standing for itself and
% homogeneity dropped (law of homogeneity explained once); powers in words →
% exponents; « in » → product; « æquatur », a brace, or a vertical stroke → =;
% cossic N, Q, C, QQ, QC, CC → x, x², …, x⁶; L, LC, LQQ, LQC, LV, LVC →
% radicals; the double stroke « = » → the signed difference (an unspecified
% order), or the margin's « + » where the margin corrects it; positivity of
% roots stated as the reading's standing convention; astronomy: prostaphaeresis
% → equation, double elongation → 2η, sines and « Sinus Totus » → sin, tan.
%
% Modern names attached, each where the statement coincides, later work
% footnoted as later: depressed equation (removal of the second term);
% Vieta's substitution; Cardano's formula; casus irreducibilis and the
% trigonometric solution of the cubic; the discriminant of x³ − px − q
% (4p³ − 27q² = (A−B)²(2A+B)²(A+2B)², the reading's factorisation of the
% second hand's inequality); AM–GM; Ferrari's method and the resolvent cubic;
% Vieta's formulas / elementary symmetric functions (Cap. XVI of the first
% treatise; Cap. XIIII of the second); divisibility of xⁿ ± yⁿ by x ± y;
% extreme and mean ratio (golden ratio), and, in a note only, the n-bonacci
% constants; the factor theorem; Descartes's rule of signs and the rational
% root theorem (notes only, as what the text tacitly needs); Tschirnhaus
% (note only: the leaves carry the linear case). Astronomy: equation (of
% centre), eccentric and epicycle, arcsin(r/R) for the maximum equation,
% variation (first prostaphaeresis, ∝ sin 2η), evection (Copernican device of
% the varying epicycle; the « anomalia mixta » 2η − M of the restored
% hypothesis), and, to first order, the three terms of lunar longitude with
% amplitudes 6°14′, 1°15′, 37′ against the modern 6°17′, 1°16′, 39′30″; in a
% note, the first-order equivalence of the planetary « adlateralis »
% construction (2φ sin σ, φ cos σ) with the equant and Kepler motion.
%
% Every computation the leaves offer was redone; the recomputations are this
% reading's own. Departures from the leaf, each footnoted where it occurs:
% leaf A, case 1 of « Facillima resolutio » (the written reduced equation
% belongs to case 2; true: a quadratic in E³ / in A²); leaf A, the climactic
% construction (F = D⁵/Z⁴ as read, D⁶/Z⁵ needed); Cap. V the series 27,18,12,8
% read 8,12,18,27, « B in Z » for D, « secundæ » for primæ; Cap. VI the
% condition phrase with « medietatis » (factor 4), the equality case D² = 4B³,
% « faciunt B planum » for the coefficient 3B; the porism (v. 23) restated as
% « zero or two positive roots »; Cap. XIII th. X–XII (v. 41–43) the division bar
% over x⁴, « Z plano plano » for Z³; Cap. XV « S² − B² » for B² − S²; the
% syncrisis signs (« = » for +, « A gradu + B gradu » for + E^k, « D » for B,
% « per A » for A − E); Cap. XVIII « 124,960 » for 249 920, « 16/15 » for
% 16/5, « B plano » for planoplano, the roots 25, 100 given for x²; Cap. XIX
% « 4,618,812 » for 5 153 632; Cap. XX the exponent of Z (quadrato for
% quadratocubo) and « planoplano »; Emendatio: the uncorrected « 30Q » (v. 78),
% « 12500 », « 3 vel 6 », « −46Q/160 », the ratios 8 and 18/2, « 8Q », and
% every other margin correction received; the table « 46 bis » row IIII sign;
% the posterior-formula theorem IIII « +4ZG » and « 64N »; Cap. IX prop. II
% « 1Q − 3N » for 1Q + 3N; prop. III as first written (the margin's form
% followed) and its example « L2 » for √8; prop. IX « +10Q » for −10Q; prop. X
% « L3 − 1 » for √3 + 1; Cap. X prop. VI, only the minus sign gives a root;
% the margin's « eadem sequuntur » for 3Bx − x³ = 2Z (no positive root under
% Z² > B³); Cap. XI prop. V « B in A qq » for B²x⁴; astronomy: 9 150 131 for
% ≈ 9 150 410 (and « EF » for AE), 0°32 15/29′ kept for 941 (≈ 0°32⅓′),
% 2 325 665 for 666, « complementum AEG » for EAG, « IQ dupla ad IH »
% inverted, the contraction quadrants, « talium 1 qualium GE 10 » against the
% numbers 1 : 5. Steps supplied: roots the page omits (18 at v. 9, 26 at
% v. 11; second positive roots at v. 10, 13, 74, 88 (1.4026) ; the extra root
% B of Cap. XIIII th. I and the roots lost by division in Emendatio Cap. IX
% props. II and IX); the justification of the irrationality test (rational
% root theorem); the triangular unciae real only for a negated affection; the
% constitution of Cap. XVIII « De reliquis » (B = u³ + v³ + a v²,
% Z = u²v²(a + b + c)); the discriminant factorisation; the modern amplitudes.
% Left unreconstructed: the theorem of leaf B; the margin example
% « 1C + 9Q − 24N = 33 » (v. 32, no rational root as read); the margins of
% v. 42, 55, 137–138 (partly), 153, 156; the squaring passage of Cap. X
% prop. VI (v. 134); the Mercury « factæ » (v. 163).
\documentclass[11pt,a4paper]{article}
% The preamble shared by every transcription of the mathematicians' manuscripts in Gallica.
%
% It rests on one idea: **uncertainty compounds, it does not resolve.** A
% transcription that smooths over a doubtful word has destroyed the only thing
% it was made for — a reader must be able to tell, at every point, what was on
% the page and what was supplied. The apparatus macros below are therefore the
% critical apparatus, not decoration.
%
% The same commands are recognised by `scripts/render.mjs`, which produces the
% site's reading view, and by `scripts/tei.mjs`. A macro added here must be
% added there too, or rendering fails loudly — which is the wanted behaviour.
%
% Comments here are English, like the rest of the repository. The documents
% this preamble sets are French, and that is deliberate: see the skills.

\ProvidesPackage{germain}[2026/09/15 Transcriptions of mathematicians' manuscripts in Gallica]

\usepackage{amsmath,amssymb,amsthm}
\usepackage{mathrsfs}
% Kept from the parent preamble so the renderer's diagram support stays
% available; these manuscripts draw no commutative diagrams, and nothing here uses it.
\usepackage{tikz-cd}
\usepackage[english,french]{babel}
\usepackage{fontspec}

% --- Greek ------------------------------------------------------------------
% The text face stays fontspec's default, Latin Modern, which has no Greek:
% XeTeX drops every glyph it lacks with a line in the log and nothing on the
% page, and the Greek words the manuscripts quote vanished from the PDFs. So
% the Greek and Greek Extended blocks get a XeTeX character class of their own,
% and entering or leaving a run of it switches the family to CMU Serif — the
% same Computer Modern design, with polytonic Greek — and back. Only the
% family changes: a Greek word in \textit or \textbf keeps its shape and series.
% The transitions to and from every other class carry the tokens that class
% has towards an ordinary letter, so French spacing before « ; » or inside
% « » still holds after a Greek word. No grouping, so no brace or \endgroup
% can come out unbalanced; maths is left alone, since XeTeX inserts nothing
% there. Kept identical in `exercices.sty`.
\makeatletter
\newfontfamily\germainGreekfont{cmun}[
  Extension=.otf, NFSSFamily=germaingreek,
  UprightFont=*rm, ItalicFont=*ti, SlantedFont=*sl,
  BoldFont=*bx, BoldItalicFont=*bi, BoldSlantedFont=*bl]
\newXeTeXintercharclass\germain@greek
\def\germain@greekrange#1#2{%
  \@tempcnta=#1\relax
  \loop
    \XeTeXcharclass\@tempcnta=\germain@greek
    \ifnum\@tempcnta<#2\advance\@tempcnta\@ne\repeat}
\germain@greekrange{"0370}{"03FF}
\germain@greekrange{"1F00}{"1FFF}
\def\germain@greekprev{\rmdefault}
\def\germain@greekin{%
  \edef\germain@greekprev{\f@family}\fontfamily{germaingreek}\selectfont}
\def\germain@greekout{\fontfamily{\germain@greekprev}\selectfont}
\def\germain@greekjoin{%
  \XeTeXinterchartokenstate=\@ne
  \@tempcnta=\z@
  \loop
    \ifnum\@tempcnta=\germain@greek\else
      \XeTeXinterchartoks\germain@greek\@tempcnta=\expandafter{%
        \expandafter\germain@greekout\the\XeTeXinterchartoks\z@\@tempcnta}%
      \XeTeXinterchartoks\@tempcnta\germain@greek=\expandafter{%
        \the\XeTeXinterchartoks\@tempcnta\z@\germain@greekin}%
    \fi
    \ifnum\@tempcnta<4095\advance\@tempcnta\@ne\repeat}
% babel-french sets its own classes, and saves and restores the interchar
% state, each time a language is selected: join again after it does.
\addto\extrasfrench{\germain@greekjoin}
\addto\extrasenglish{\germain@greekjoin}
\AtBeginDocument{\germain@greekjoin}
\makeatother

\usepackage{geometry}
\geometry{a4paper,margin=25mm,marginparwidth=28mm}
\usepackage{marginnote}
\usepackage{xcolor}
\usepackage[normalem]{ulem}
\setlength{\emergencystretch}{3em}
\usepackage{hyperref}
\hypersetup{colorlinks=true,linkcolor=black,urlcolor=black,pdfborder={0 0 0}}

% --- Batch metadata ---------------------------------------------------------
% Taken as they stand by the rendering script, which shows them at the head of
% the reading view. They fix what the file claims to cover. The macro names are
% the parent project's, so that the scripts stay shared; \folder here means the
% volume's slug — `fr-9115` — and \pages counts Gallica views.

\newcommand{\folder}[1]{\gdef\grothFolder{#1}}
\newcommand{\batch}[1]{\gdef\grothBatch{#1}}
\newcommand{\pages}[2]{\gdef\grothFirst{#1}\gdef\grothLast{#2}}
\newcommand{\foldertitle}[1]{\gdef\grothFoldertitle{#1}}

% The holder's own shelfmark, printed as they write it: « Français 9115 ».
\newcommand{\shelfmark}[1]{\gdef\grothShelfmark{#1}}

% The dating, copied verbatim from the holder's catalogue — never inferred
% here. « XIXe siècle » is the BnF's; a pass that dates a leaf from its content
% says so in a \note{}, as its own inference.
\newcommand{\dating}[1]{\gdef\grothDating{#1}}

% The status notice, printed in the title block and repeated as a diagonal
% watermark on every page of the PDF. The manuscripts are in the public
% domain; what this line declares is not a right but a quality — an
% unverified first pass by a machine. The skills write the line; a file
% without it gets no watermark, so leaving it out is visible in the output.
\newcommand{\watermark}[1]{\gdef\grothWatermark{#1}}

\gdef\grothFolder{?}\gdef\grothBatch{}\gdef\grothFirst{?}\gdef\grothLast{?}%
\gdef\grothFoldertitle{}\gdef\grothShelfmark{}\gdef\grothDating{}\gdef\grothWatermark{}

% --- The résumé -------------------------------------------------------------
\newenvironment{resume}
  {\small\setlength{\parskip}{0.45em}}
  {\par\medskip\hrule\medskip}

% --- Keywords ---------------------------------------------------------------
% In English, closing the modernised reading's résumé: the modern vocabulary
% under which to search for what the leaves construct. The manifest extracts
% them to tag the volume — this line is the single source of the tags.
\newcommand{\keywords}[1]{\par\smallskip\noindent
  {\footnotesize\textsc{Keywords} --- \textit{#1}}\par}

% --- The critical apparatus -------------------------------------------------

% The Gallica view number, in the margin. This is the anchor that lets a
% disagreement over a formula be settled by looking at *one* image rather than
% a volume of 750. The site uses it to turn the facsimile as the transcription
% is scrolled. A view is one scanned image: usually one side of a leaf.
\newcommand{\page}[1]{%
  \marginnote{\footnotesize\color{gray!70}\textsf{v.\,#1}}%
  \label{page:#1}\ignorespaces}

% The BnF's pencilled foliation, where the leaf carries it and a pass can read
% it: « 198r ». This is what the literature cites — « ff. 198r–208v » — and
% the one line that turns a citation into a view. Recorded, never inferred:
% a folio a pass cannot read is not written.
\newcommand{\folio}[1]{%
  \marginnote{\footnotesize\color{gray!50}\textsf{f.\,#1}}\ignorespaces}

% The modernised reading's coarser counterpart to \page{}: which views a
% section is drawn from.
\newcommand{\pagerange}[2]{%
  \marginnote{\footnotesize\color{gray!70}\textsf{v.\,#1--#2}}%
  \ignorespaces}

% Illegible: never guessed.
\newcommand{\ill}{\textcolor{red!60!black}{[\,\ldots\,]}}

% A doubtful reading: the word is given, and flagged as uncertain.
\newcommand{\uncertain}[1]{\uline{#1}}

% An editorial addition — a missing letter, an implied word.
\newcommand{\add}[1]{\textcolor{blue!50!black}{[#1]}}

% Struck out by the writer. What was crossed out is part of the document.
\newcommand{\struck}[1]{\sout{#1}}

% The transcriber's note, on the state of the page or the mathematics.
\newcommand{\note}[1]{\par\smallskip{\small\color{gray!60!black}\textit{[#1]}}\par\smallskip}

% A marginal note of hers, distinct from the body of the text.
\newcommand{\marginal}[1]{\par\smallskip\noindent\hspace{1em}%
  {\small\color{gray!60!black}#1}\par\smallskip}

% --- Title ------------------------------------------------------------------
% The title block is French, like the documents themselves.

\AddToHook{shipout/background}{%
  \ifx\grothWatermark\empty\else
    \begin{tikzpicture}[remember picture, overlay]
      \node[rotate=52, scale=3.4, align=center, text=gray!18,
            font=\sffamily\bfseries]
        at (current page.center) {\grothWatermark};
    \end{tikzpicture}%
  \fi}

\AtBeginDocument{%
  \begin{center}
    {\large\bfseries Manuscrits de mathématiciens (Gallica) ---
      \ifx\grothShelfmark\empty\grothFolder\else\grothShelfmark\fi}\\[2pt]
    {\itshape\grothFoldertitle}\\[4pt]
    {\small vues \grothFirst--\grothLast
      \ifx\grothBatch\empty\else\ (lot \grothBatch)\fi}\\[2pt]
    \ifx\grothDating\empty\else
      {\small\color{gray!60!black}Datation du catalogue : \grothDating}\\[2pt]
    \fi
    \ifx\grothWatermark\empty\else
      {\small\bfseries\color{red!45!gray}\grothWatermark}\\[2pt]
    \fi
    {\footnotesize\color{gray!60!black}Transcription établie d'après les images IIIF de Gallica.
     Source gallica.bnf.fr / Bibliothèque nationale de France.\\
     Les lectures incertaines sont soulignées, les illisibles notées [\,\ldots\,],
     les ajouts entre crochets ; \textsf{v.} numérote les vues de Gallica,
     \textsf{f.} la foliotation au crayon de la BnF.}
  \end{center}
  \vspace{1em}}

\endinput


\folder{nal-1644}
\shelfmark{NAL 1644}
\pages{1}{163}
\dating{}
\watermark{Lecture automatique\\interprétation non vérifiée}
\foldertitle{Viète (François). Papiers divers. NAL 1644 — lecture modernisée du volume entier}

\begin{document}

\section*{Résumé}

\begin{resume}
Ce volume réunit trois pièces sous une couverture qui porte « Francisci Vietæ
opera nonnulla », quelques œuvres de François Viète. Les deux premières sont
deux traités d'algèbre en latin, d'une même écriture selon les transcriptions, titrés sur leurs
feuillets au nom de Viète : « De Recognitione Æquationum » (vues 5 à 68), sur la
manière de reconnaître comment une équation est faite, et « De æquationum
Emendatione tractatus secundus » (vues 69 à 141), sur la manière de la préparer
pour la résoudre. Qui les a écrits, les feuillets ne le disent pas ; trois
notes, toutes de lecture douteuse, pourraient en faire la copie d'un
« exemplar ». La troisième pièce (vues 142 à 163) est un dossier d'astronomie
sans nom : les tables des inégalités de la Lune, une « restitution » de
l'hypothèse lunaire de Tycho Brahe, le début d'un chapitre qui reproche à
Copernic sa mauvaise géométrie, et des fragments sur les planètes, qui
s'arrêtent au milieu d'une phrase sur Mercure. Aucun feuillet n'est daté.

Le problème des deux traités est celui de toute l'algèbre de leur temps :
résoudre les équations du troisième et du quatrième degré, et, pour les
autres, en calculer numériquement les racines. Le premier traité dit
d'emblée ce qu'il apporte : la doctrine générale de la résolution numérique,
il la « forme et l'achève », parce que « les équations ont souvent besoin
d'une préparation avant de se laisser facilement résoudre ». Pour préparer une
équation, il faut savoir de quoi elle est faite : de quelles racines, par quel
changement d'inconnue à partir d'une équation plus simple, avec quels
coefficients. Tout est écrit dans une algèbre de lettres, où chaque grandeur
porte sa dimension, et où les racines sont toujours positives.

L'idée maîtresse est le changement d'inconnue. Translater la racine d'une
fraction du second coefficient fait disparaître le terme qui suit la plus haute
puissance ; prendre pour inconnue le quotient du terme connu par la racine
change les signes ; poser $x = (B - y^2)/y$ ramène la cubique $x^3 + 3Bx = 2Z$ à
une équation du second degré en $y^3$, d'où la formule de Cardan ; ajouter aux
deux membres d'une quartique ce qui en fait deux carrés ramène la quartique à
une cubique auxiliaire, puis à deux quadratiques. Une seconde idée, la
« syncrise », compare deux équations de même forme et de mêmes coefficients qui
ont des racines différentes : en les soustrayant, on exprime les coefficients
par les racines — le coefficient est leur somme, le terme connu leur produit —,
et le dernier chapitre du second traité l'énonce pour deux, trois, quatre et
cinq racines. Le cas où la cubique a trois racines réelles, où la formule de
Cardan échoue, est résolu par la trisection de l'angle. Les dizaines d'exemples
chiffrés sont presque tous justes ; les fautes que l'écrivain n'a pas lui-même
corrigées en marge sont signalées, et quelques énoncés trop forts sont rectifiés.
Le dossier d'astronomie décompose le mouvement de la Lune en inégalités qui
dépendent de l'anomalie, du double de l'élongation au Soleil et de leur
différence ; l'écrivain juge que Tycho « n'est nullement géomètre » et écrit :
« Ego suam hypothesim ita restituo », je restitue son hypothèse ainsi.

Ce que ces feuillets construisent porte aujourd'hui des noms précis : les
formules de Viète et les fonctions symétriques élémentaires des racines ;
l'équation réduite, privée de son second terme ; la substitution de Viète et la
formule de Cardan ; le cas irréductible et la résolution trigonométrique de la
cubique ; la méthode de Ferrari et la cubique résolvante ; le discriminant ; la
moyenne et extrême raison, c'est-à-dire le nombre d'or. En astronomie :
l'équation du centre, l'évection et la variation, les trois grandes inégalités
de la longitude de la Lune, et les modèles à excentrique et épicycle.

\keywords{theory of equations, law of homogeneity, Vieta's formulas, elementary symmetric functions, transformation of equations, depressed equation, Vieta's substitution, cubic equations, Cardano's formula, casus irreducibilis, trigonometric solution of the cubic, angle trisection, discriminant, inequality of arithmetic and geometric means, quartic equations, Ferrari's method, resolvent cubic, geometric progressions, extreme and mean ratio, golden ratio, lunar theory, equation of centre, evection, lunar variation, epicycle and eccentric, Tycho Brahe's lunar theory, Copernicus, planetary theory}
\end{resume}

\section*{Ce que contient le volume}

\pagerange{1}{163}
Cent soixante-trois vues, une page par vue, feuillets de tailles inégales posés
les uns sur les autres, de sorte que le haut ou le pied d'un feuillet voisin
dépasse sur la plupart des images. Sous une couverture de papier fort qui porte,
d'une grande écriture ancienne, « Francisci Vietæ opera nonnulla »\footnote{« Quelques
œuvres de François Viète ». La transcription de la vue 1 relève aussi, d'une
main plus petite et postérieure, « imprimés, pag. 84 de l'edition d'Elzevir
1646. in folio p 158 », puis une petite table : « de æquationum recognitione 1. /
De emendatione æquationum 32 » (« pag. 84 », « p » et « de » lus avec doute,
« emendatione » de même). Les nombres 1 et 32 de cette table sont ceux que
porte, en tête de page, la série de chiffres à l'encre décrite plus bas, aux
vues 5 et 69, là où commencent les deux traités : c'est un rapprochement fait
ici, non une lecture de la table. Les renvois à l'édition de 1646 sont d'un
lecteur ; ils n'ont été confrontés à aucun imprimé.}, le volume réunit trois
pièces, précédées de trois petits feuillets lettrés.

\begin{itemize}
  \item \textbf{Vues 2 à 4 — trois petits feuillets lettrés A, B, C}, d'une
    écriture serrée différente de celle des traités : deux réductions
    d'équations du sixième degré et un théorème sur les suites de
    proportionnelles (A, B) ; le feuillet C est blanc.
  \item \textbf{Vues 5 à 68 — « Francisci Vietæ De Recognitione Æquationum
    tractatus. »}, ainsi titré en tête de la vue 5, chapitres 1 à XXI : comment
    une équation est faite à partir de ses racines, et comment, la voyant, on
    reconnaît de quoi elle est faite.
  \item \textbf{Vues 69 à 141 — « Francisci Vietæ. De æquationum Emendatione
    tractatus secundus. »}, ainsi titré en tête de la vue 69, chapitres I à
    XIIII, qui s'achève sur « Finis » : comment préparer une équation pour la
    résoudre.
  \item \textbf{Vues 142 à 163 — des feuillets d'astronomie, sans nom}, d'une
    autre matière : les prostaphérèses de la Lune et le cadre de leurs tables,
    l'hypothèse lunaire de Tycho Brahe « restituée » par l'écrivain, un
    chapitre XII intitulé « Harmonici Ptolemaici Constructio » qui s'ouvre sur
    une charge contre la géométrie de Copernic, et deux fragments de théorie
    planétaire, dont le dernier porte sur Mercure et s'interrompt au milieu
    d'une phrase.
\end{itemize}

Les deux traités se suivent sans lacune : le second commence en tête du
feuillet qui suit la dernière page du premier, et il est dit « secundus ». Leur
écriture est, selon les transcriptions, une même cursive régulière, qui se
corrige elle-même dans la marge — un mot souligné dans la ligne, la bonne
leçon soulignée en regard. D'autres mains interviennent : une main rapide à
l'encre noire, dans les marges des vues 13 à 18, qui démontre ce que le texte
énonce et parle de ses propres travaux ; une petite main pâle à la vue 20 ; une
écriture plus serrée dans les marges des vues 61 à 66 et une couche de
corrections d'encre plus noire aux vues 121 à 140, dont les transcriptions ne
décident pas si elles sont de l'écrivain ou d'un autre ; enfin un lecteur plus
tardif, qui renvoie aux pages d'une édition in-folio de 1646 (vues 1, 5, 69 et
141).\footnote{Les renvois sont : « page 84 de l'edition d'Elzevir 1646 in
folio. » (vue 5, en tête du premier traité) ; « p. 127. editionis 1646. »
(vue 69, en tête du second) ; « pag. 158 editionis elzevir. anni 1646,
infolio. » (vue 141, à la fin du second). Ce sont des pages d'un livre imprimé,
non du manuscrit. Aucun n'a été vérifié : cette lecture n'a ouvert aucune
édition imprimée des œuvres de Viète, ni pour contrôler ces renvois, ni pour
combler une lacune, ni pour trancher une leçon.}

Le nom qu'on lit sur ces feuillets est celui des deux titres, écrits de la main
du texte, et celui de la couverture. Qui a tenu la plume, les transcriptions ne
le disent pas, et cette lecture ne le dit pas davantage. Trois notes, toutes
douteuses, pourraient faire des traités une copie : « transcripta ex
exemplari », en tête d'une addition marginale encadrée de la vue 100
(« copiée sur l'exemplaire », lecture probable, non certaine) ; « deerant hæc
in exemplari », écrit entre les lignes puis biffé à la vue 114 (« ceci manquait
dans l'exemplaire », lecture très incertaine) ; et, à la vue 20, une note d'une
petite main pâle appelée par une croix, « in exemplari \emph{excripto} / vel
potius \emph{adscriptum} », qui n'est lue qu'en partie.\footnote{Les mots en
italique dans cette dernière note sont ceux que la transcription donne comme
incertains. Si ces trois lectures sont bonnes, un modèle (« exemplar ») a été
recopié ; elles ne disent ni quel modèle, ni par qui, ni quand. Cette lecture
les donne pour ce qu'elles sont : des indices douteux.} Les feuillets
d'astronomie ne portent aucun nom d'auteur, et cette lecture ne les attribue à
personne. Aucun feuillet n'est daté ; le catalogue ne donne pas de datation,
et il n'en est inféré aucune.

Une seule série de chiffres court sur tout le volume : un chiffre à l'encre en
haut à droite de chaque recto, de 1 (vue 5) à 79 (vue 162), sans lacune, avec
« 5 bis », « 16 bis », « 46 bis » et « 47 bis » pour des béquets insérés. Elle
traverse les trois pièces, et l'écrivain la cite lui-même : « ad sinistram
paginam fol. 18. » à la vue 43, « vide annotata pag. 48 » à la vue 100. C'est
donc le compte des feuillets de celui qui a écrit ou copié, non la foliotation
de la bibliothèque ; les transcriptions ne l'ont pas écrite en \emph{folio}, et
cette lecture ne cite que des vues.\footnote{Le « 16 bis » de la vue 38 est
d'une autre encre et d'une autre main que la fiche qu'il numérote, laquelle
porte, de la main du texte, « ad faciem aversam foly 16. » (« au revers du
feuillet 16 »). Les autres chiffres des feuillets sont les cotes et cachets
(« Libri 1201 », « Lat. nouv. acq. 1644 », « R.c. 8070 (73) », « ACQ. 8070
(LIBRI) », vues 1, 2 et 5), un « 8 » au crayon sur la couverture, et, dans le
texte, les numéros de chapitres, de théorèmes et de propositions, souvent
corrigés par l'écrivain (« Cap. VIII » récrit « IX » à la vue 23, « XIIII »
récrit « XV » à la vue 45, « XX » ramené à « X » dans la marge à la vue 130,
etc.). Cette lecture suit la numérotation corrigée.}

Parce que les deux traités se tiennent et que l'astronomie en est séparée, la
lecture qui suit leur donne deux suites de sections : l'algèbre d'abord,
chapitre par chapitre dans l'ordre des feuillets, puis l'astronomie. Elle ne
cherche pas de fil entre les deux.

\section*{Les trois petits feuillets A, B et C}

\pagerange{2}{4}
Trois feuillets plus petits que la page, lettrés A, B et C en haut à droite,
d'une écriture compacte que la transcription dit différente de celle des
traités, et que rien ne rattache à eux.\footnote{La marge du feuillet A porte,
d'une encre plus noire et d'une autre main, « Vid. fol. 2 ». Le feuillet C
(vue 4) est blanc et n'est pas transcrit.} Le feuillet A traite d'équations du
sixième degré « a radice plana solidave », c'est-à-dire d'équations qui ne
contiennent l'inconnue $E$ que par $E^3$ ou par $E^2$.

\subsection*{La « duplicata hypostasis »}

L'équation
\[ E^6 + Z^3 E^3 = D^6 \]
est ramenée à la géométrie des moyennes proportionnelles.\footnote{Le feuillet
écrit « E cc + Z s. in E c. æq. D ss. » : $E$ cubo-cubus, plus $Z$ solide
multiplié par $E$ cube, égale $D$ solido-solide. Sur la notation, voir la
section qui suit.} On pose $B = Z^3/D^2$ ; on cherche la plus petite, $F$, de
trois proportionnelles de moyenne $D$ et de différence des extrêmes $B$,
c'est-à-dire $F^2 + BF = D^2$ ; on fait $E^3 = D^2 F$ ; et $E$ est la première
des deux moyennes proportionnelles entre $D$ et $F$. La construction est
juste : si $u = D^2 F$, alors
\[ u^2 + Z^3 u = D^4\bigl(F^2 + BF\bigr) = D^6 , \]
et la première moyenne entre $D$ et $F$ est bien $\sqrt[3]{D^2F}$. C'est la
résolution d'une équation du second degré en $E^3$, suivie d'une extraction de
racine cubique, faite en géométrie.

\subsection*{La « climactica paraplerosis »}

L'équation $E^6 + Z^4E^2 = D^6$ est, de même, une équation du troisième degré en
$E^2$. Le feuillet la ramène à la construction de quatre proportionnelles
$Z, s, s^2/Z, s^3/Z^2$ dont la somme du second et du quatrième terme est une
longueur donnée $F$, puis fait $E^2 = Zs$. La réduction est juste à condition
que
\[ s + \frac{s^3}{Z^2} = \frac{D^6}{Z^5}, \]
puisqu'alors $E^6 + Z^4E^2 = Z^5\bigl(s + s^3/Z^2\bigr) = D^6$. Telle que la
transcription la lit, la construction de $F$ donne $F = D^5/Z^4$ et non
$D^6/Z^5$ : il manque un facteur $D/Z$.\footnote{Le feuillet pose la série
$Z, D, D^2/Z, D^3/Z^2$, puis « ut $Z$ ad quartam ita $D\,q.$ ad $C\,q.$ » et « ut
$Z\,p.$ ad $C$ plano ita $Z$ ad $F$ », d'où $C^2 = D^5/Z^3$ et $F = D^5/Z^4$. Une
proportion est donc mal lue ou omise ; la transcription note que la lettre lue
$C$ pourrait être un $G$, et qu'une première proportion, « fiat ut $Z\,ppl.$ ad
$D\,pp.$ », reste en suspens devant la série. Cette lecture ne restitue pas la
proportion manquante.}

\subsection*{« Facillima resolutio æquationum cubicarum affectarum »}

Le paragraphe suivant annonce que ces équations « se ramènent aux quadratiques
et à la racine simple », et traite deux cas. Pour
\[ E^6 + Z^4E^2 = B^6 , \]
la substitution $E^2 = ZA$ donne bien
\[ A^3 + Z^2 A = \frac{B^6}{Z^3}, \]
une cubique en la racine simple $A$, comme l'écrit le feuillet. Pour le premier
cas, $E^6 \pm Z^3E^3 = B^6$, le feuillet pose $E^3/Z = A^2$ et écrit la même
équation $A^3 + Z^2A = B^6/Z^3$ ; ce n'est pas ce que donne la substitution.
Avec $E^3 = ZA^2$ on obtient
\[ A^4 \pm Z^2A^2 = \frac{B^6}{Z^2}, \]
une équation du second degré en $A^2$ ; et, plus simplement, l'équation
proposée est du second degré en $E^3$.\footnote{Le feuillet écrit, pour le
premier cas, « fingatur $\frac{E\,c.}{Z}$ æquari $A\,q.$ » puis « æquatio namq;
orta sic se habebit $A\,c. + Z\,q.$ in $A$ æq. $\frac{B\,ss.}{Z\,c.}$ », c'est-à-dire
l'équation qui convient au second cas. Il ne donne que le signe $+$ dans les
deux équations réduites, là où l'énoncé porte « plus minusve ». La finale de
« $B\,ss.$ » compte un jambage de trop dans le premier cas, et le mot après
« simplici » est biffé (« quadr\ldots » ?).}

\subsection*{Le théorème du feuillet B}

Le feuillet B énonce un théorème sur les termes « pris alternativement » d'une
suite de continues proportionnelles, en nombre pair puis en nombre impair, et
renvoie sa preuve à la « méthode analytique ». Tel que la transcription le lit,
avec deux additions interlinéaires, l'énoncé ne se laisse pas réduire à une
identité qu'on puisse vérifier sans deviner la construction de la phrase ; cette
lecture ne le reformule donc pas.\footnote{Texte lu : « Si fuerit series
continue proportionalium numero parium, differentia postremarum cum [alterne
sumptarum] accepta prima, minus differentia cum alterne sumptarum æquatur
extremæ minori plus eadem maiori », puis le cas impair ; « Debet autem in serie
proposita intelligi prima minor. quod methodo Analyticâ facile patebit. » Le mot
lu « cum » est partout abrégé. Les chapitres XVII et XIX du premier traité
contiennent des énoncés voisins sur les différences de puissances de termes
alternés ; aucun ne coïncide avec celui-ci tel qu'il est lu.}

\section*{La notation des feuillets, et sa traduction}

\pagerange{5}{23}
Les traités écrivent l'algèbre en « espèces » : des lettres pour les grandeurs,
des mots pour les opérations. On la traduit ici une fois pour toutes ; la
forme du feuillet est rappelée en note là où elle importe.

\begin{itemize}
  \item \textbf{L'inconnue} est $A$ ; une nouvelle inconnue, après
    transformation, est $E$. On les écrit $x$ et $y$. Les grandeurs données sont
    des consonnes, $B$, $D$, $G$, $H$, $S$, $X$, $Z$\ldots, que l'on garde.
  \item \textbf{L'homogénéité.} Chaque grandeur donnée porte sa dimension :
    « $B$ planum » est une aire, « $D$ solidum » un volume, « $Z$ plano-planum »
    une grandeur de dimension 4, et tous les termes d'une équation ont la même
    dimension. Quand le feuillet nomme ainsi une grandeur par sa dimension, la
    lettre désigne ici cette grandeur elle-même, un nombre ; quand il écrit une
    puissance (« $B$ quadratum », « $B$ cubus »), on écrit la puissance, $B^2$,
    $B^3$. Ainsi « $A$ cubus plus $B$ plano ter in $A$ æquatur $D$ solido »
    devient $x^3 + 3Bx = D$. La loi d'homogénéité est une contrainte que
    l'algèbre d'aujourd'hui n'impose pas ; la supprimer ne change aucun
    résultat, puisqu'un choix d'unité rend toutes les grandeurs numériques.
  \item \textbf{Les puissances} sont en mots : quadratum, cubus,
    quadrato-quadratum, quadrato-cubus, cubo-cubus, c'est-à-dire $x^2$ à $x^6$ ;
    « latus » est la racine. « In » est le produit ; « bis », « ter »,
    « quater » sont les coefficients 2, 3, 4. Il n'y a pas de signe d'égalité :
    « æquatur », « æquabitur », ou une accolade entre les deux membres.
  \item \textbf{Les exemples numériques} sont en signes cossiques, coefficient
    d'abord : N le nombre (la racine), Q le carré, C le cube, QQ, QC, CC les
    puissances quatrième, cinquième et sixième. « $1Q + 10N$ æquatur 144 » est
    $x^2 + 10x = 144$, et « fit $1N$ 8 » donne la racine. « $L\,144$ » est
    $\sqrt{144}$ ; $LC$, $LQQ$, $LQC$ sont les racines cubique, quatrième et
    cinquième ; « $LV$ » (latus universale) est la racine de toute la somme qui
    suit.
  \item \textbf{Le double trait « = ».} Là où le calcul demande une
    soustraction, les feuillets portent souvent un double trait horizontal, que
    la main distingue du moins simple. C'est le signe d'une différence dont
    l'ordre n'est pas donné : « $Z$ plano $=$ $B$ quadrato » veut dire « la
    différence entre $Z$ et $B^2$, quelle que soit la plus grande ». Comme
    l'algèbre d'aujourd'hui admet les nombres négatifs, on l'écrit ici par la
    différence signée, $Z - B^2$, qui couvre les deux cas ; là où la marge le
    corrige en « $+$ », on suit la correction et on le dit.
\end{itemize}

Une convention de fond accompagne cette notation, et il faut la dire parce que
les énoncés en dépendent : \textbf{les racines des feuillets sont positives}.
Une équation n'y a de racine que positive ; c'est pourquoi chaque forme
d'équation est écrite avec tous ses termes positifs de part et d'autre, et
pourquoi $x^2 + bx = c$, $x^2 - bx = c$ et $bx - x^2 = c$ sont trois espèces
différentes, que les chapitres traitent séparément. Chaque fois qu'un énoncé
ci-dessous dit « la racine », il faut entendre « la racine positive », et
chaque fois qu'une équation a deux racines positives, le traité le dit ou on le
dit pour lui.\footnote{Les racines négatives, que le traité ne nomme jamais,
sont données en note quand elles éclairent un énoncé ; ce sont des ajouts de
cette lecture.}

\section*{Le « De Recognitione » : ce qu'il se propose}

\pagerange{5}{6}
Le traité s'ouvre en disant à quoi il sert :

\begin{quote}
« Generalem et generaliter traditam de numerosa potestatum resolutione
doctrinam informat ac perficit tractatus de recognitione æquationum.
Præparatione enim indigent æquationes sæpenumero antequam faciliter
explicentur. »\footnote{« La doctrine générale, et transmise en général, de la
résolution numérique des puissances, le traité de la reconnaissance des
équations lui donne sa forme et l'achève. Car les équations ont souvent besoin
d'une préparation avant de se laisser facilement résoudre. » « Explicare » une
équation, dans ces traités, c'est la résoudre ; « potestas » est la puissance la
plus haute de l'inconnue, et, par extension, l'équation qu'elle domine.}
\end{quote}

Préparer une équation, c'est la transformer ; mais pour savoir comment, il faut
d'abord savoir comment elle est faite — sa « constitution ». Celle-ci, dit le
chapitre 1, se reconnaît par trois voies : la \emph{zétèse}, qui retrouve
l'équation comme la mise en équation d'un problème de proportions ; le
\emph{plasme}, qui la voit naître d'une équation plus simple par un changement
d'inconnue ; et la \emph{syncrise}, qui compare deux équations de même forme et
de mêmes données.\footnote{« Æqualitatum constitutio potissimo deprehenditur,
Zetesi, plasmate, et Syncrisi » (vue 5). Le premier paragraphe du chapitre dit
aussi quand la préparation s'impose : quand les termes négatifs l'emportent,
quand il y a des fractions ou des irrationnelles. La phrase du second
paragraphe, qui oppose le géomètre, pour qui fractions et irrationnelles ne
font pas obstacle, à l'arithméticien, contient trois mots grecs dont deux ne
sont lus qu'en partie (« πολυ\ldots » (polu\ldots), « ἀρρηξία ? ἢ ἀλογία » (arrêxia ? ê alogia)) ; plusieurs mots
de la page sont lus avec doute (« ari\ldots{} », « solet », « offerre »,
« Ergo », « imitetur », « suscepto », « ortus », « deductionis »).} Les
chapitres suivent cet ordre.

Le chapitre II définit la similitude de deux équations : elles sont
« semblables » si elles ont même degré, les mêmes termes intermédiaires et les
mêmes signes ; si une condition de plus est requise, la similitude est
« anomale ».\footnote{« Æquationem Æquationi similem regulariter dicimus cum
utrobiq; par est potestas, seu æq. alta, ipsaq; affecta affectionibus sub pari
gradu, et eadem affectionis nota » (vue 6). La deuxième ligne du chapitre est
d'une lecture très incertaine selon la transcription.}

\section*{Les équations du second degré tirées des proportions}

\pagerange{6}{8}
Le chapitre III tire trois théorèmes d'un seul problème : trois grandeurs en
proportion continue, $a : m = m : c$, dont on connaît la moyenne $m$ et la
différence ou la somme des extrêmes.\footnote{Les trois espèces portent des noms
grecs, écrits en lettres grecques : καταφατική (kataphatikê, affirmative),
ἀποφατική (apophatikê, négative), ἀμφίβολος (amphibolos, ambiguë). Le traité appuie l'extension des rapports de
lignes à des grandeurs de toute dimension sur Campanus, « in libro de
proportione » (vue 7).} Puisque le produit des extrêmes vaut le carré de la
moyenne :

\begin{itemize}
  \item si $x^2 + bx = m^2$, $x$ est le plus petit extrême de trois
    proportionnelles de moyenne $m$ dont les extrêmes diffèrent de $b$ ;
  \item si $x^2 - bx = m^2$, $x$ est le plus grand ;
  \item si $bx - x^2 = m^2$, $x$ est l'un ou l'autre des extrêmes, dont la somme
    est $b$.
\end{itemize}

Les trois exemples portent sur la suite $8, 12, 18$ : $x^2 + 10x = 144$ donne
$8$, $x^2 - 10x = 144$ donne $18$, et $26x - x^2 = 144$ donne $8$ ou $18$. Tous
sont justes.\footnote{« $1Q + 10N.$ æquatur $144$ » : $x^2 + 10x = 144$ ; « $L.\,144$ »
est $\sqrt{144} = 12$. La troisième espèce est la seule qui ait deux racines
positives, d'où son nom d'« ambiguë ». Les deux autres ont chacune une racine
négative ($-18$ et $-8$), que le traité ne considère pas.}

\section*{Les cubiques tirées des proportions}

\pagerange{8}{13}
Les chapitres IIII et V font de même avec quatre proportionnelles
$B, x, x^2/B, x^3/B^2$. Le chapitre IIII prend pour inconnue le second terme
(les « affections sous le côté ») : puisque le cube du second vaut le carré du
premier multiplié par le quatrième,

\begin{itemize}
  \item $x^3 + B^2x = B^2Z$ quand le second et le quatrième ont pour somme $Z$ ;
  \item $x^3 - B^2x = B^2D$ quand $B$ est le plus petit extrême et que le
    quatrième surpasse le second de $D$ ;
  \item $B^2x - x^3 = B^2D$ quand $B$ est le plus grand extrême et que le second
    surpasse le quatrième de $D$.
\end{itemize}

Les exemples sont justes : $x^3 + 64x = 2496$ donne $12$ (suite $8, 12, 18,
27$) ; $x^3 - 64x = 960$ donne $12$ ; $729x - x^3 = 7290$ donne $18$ (suite
$27, 18, 12, 8$).\footnote{Le second exemple du premier théorème,
$x^3 + 729x = 18954$, ne donne pas sa racine ; c'est $18$, le second terme de la
suite $27, 18, 12, 8$ lue à rebours ($18^3 + 729 \cdot 18 = 18954$). Cette valeur
est un ajout de cette lecture.} Le traité ajoute que, dans la troisième espèce,
« le second peut être double », et le montre sur deux suites de même premier
terme $\sqrt{59319}$ et de même différence $70$ entre second et quatrième :
$\sqrt{59319}, 195, \sqrt{24375}, 125$ et $\sqrt{59319}, 78, \sqrt{624}, 8$. Ce
sont bien les deux racines positives de $59319x - x^3 = 4\,152\,330$, et
l'exemple précédent en avait deux aussi.\footnote{Vérifié : $195^2 = 38025$ et
$38025^2 = 59319 \cdot 24375$ ; $78^2 = 6084$ et $6084^2 = 59319 \cdot 624$ ;
$195^3 = 59319 \cdot 125$ et $78^3 = 59319 \cdot 8$. L'équation
$x^3 - 59319x + 4\,152\,330 = 0$ a pour racines $195$, $78$ et $-273$. De même
$729x - x^3 = 7290$ a, outre $18$, la racine positive $-9 + \sqrt{486} \approx
13{,}05$, que la page ne donne pas. Le traité montre enfin le même phénomène sur
trois proportionnelles, $36, 6, 1$ et $36, 30, 25$, dont le second surpasse le
troisième de $5$ : ce sont les deux racines de $36x - x^2 = 180$.}

Le chapitre V prend pour inconnue la somme ou la différence du premier et du
troisième terme (les « affections sous le carré ») :

\begin{itemize}
  \item si $x^3 - Bx^2 = BZ^2$, et que le second et le quatrième ont pour somme
    $Z$, $x$ est la somme du premier et du troisième ;
  \item si $x^3 + Bx^2 = BD^2$, avec $B$ le plus petit extrême et $D$ la
    différence du second et du quatrième, $x$ est la différence du premier et du
    troisième ;
  \item si $Bx^2 - x^3 = BD^2$, avec $B$ le plus grand extrême, $x$ est encore
    la différence du premier et du troisième.
\end{itemize}

La démonstration est la même dans les trois cas : si $x$ est la somme (ou la
différence) du premier et du troisième terme, le second vaut $BZ/x$, parce que
les sommes (ou les différences) de termes de même rang sont dans le rapport de
la suite ; le produit du premier par le troisième égale le carré du second, et
l'on multiplie par $x^2/B$. Les exemples sont justes : $x^3 - 8x^2 = 12168$
donne $26 = 8 + 18$ ; $x^3 + 8x^2 = 1800$ donne $10 = 18 - 8$ ;
$27x^2 - x^3 = 2700$ donne $15 = 27 - 12$.\footnote{Trois écarts du feuillet.
(1) Au premier exemple, la suite est écrite $27, 18, 12, 8$, le rang I sur $27$,
alors que le premier terme de l'exemple est $8$ : il faut la lire $8, 12, 18,
27$, dont le second et le quatrième font $12 + 27 = 39 = \sqrt{12168/8}$ ; la
racine, que la page ne donne pas, est $26$. (2) Au deuxième théorème, la fin du
calcul écrit « $B$ in $Z$ quadratum » là où l'énoncé et le calcul portent $D$.
(3) Au troisième, le zététique demande « la différence du second et du
troisième » ; l'énoncé et la suite disent « du premier et du troisième ». Enfin
$27x^2 - x^3 = 2700$ a une seconde racine positive, $6 + \sqrt{216} \approx
20{,}70$, que la page ne mentionne pas.}

\section*{Le chapitre VI : la cubique sans second terme}

\pagerange{13}{18}
Le chapitre VI est le plus important de ce premier groupe. Il donne une
« constitution singulière » de la cubique privée de son terme en $x^2$, celle
qui porte aujourd'hui la formule de Cardan, et il traite à part le cas où
cette formule échoue.

\subsection*{Deux côtés, leur produit et la différence ou la somme de leurs cubes}

\begin{itemize}
  \item Si $x^3 + 3Bx = D$, il y a deux côtés $u > v$ dont le produit est $B$ et
    dont les cubes diffèrent de $D$, et $x = u - v$.
  \item Si $x^3 - 3Bx = D$, et si $D^2 > 4B^3$, il y a deux côtés de produit $B$
    dont les cubes ont pour somme $D$, et $x = u + v$.
\end{itemize}

C'est l'identité
\[ (u \mp v)^3 \pm 3uv\,(u \mp v) = u^3 \mp v^3 . \]
Les exemples sont justes : $x^3 + 6x = 7$ donne $1 = 2 - 1$ ($2 \cdot 1 = 2$,
$8 - 1 = 7$) ; $x^3 - 6x = 9$ donne $3 = 2 + 1$. Pour trouver $u$ et $v$ au
premier cas, il suffit de résoudre une équation du second degré : $u^3$ et
$-v^3$ sont les racines de $t^2 - Dt - B^3 = 0$. On obtient ainsi, en toutes
lettres, pour $x^3 + 3Bx = D$ :
\[ x = \sqrt[3]{\sqrt{\tfrac{D^2}{4} + B^3} + \tfrac{D}{2}}
     - \sqrt[3]{\sqrt{\tfrac{D^2}{4} + B^3} - \tfrac{D}{2}} . \]
Cette formule n'est pas écrite ici ; le second traité l'écrit, au chapitre VII
de l'« Emendatio » (vues 119 à 122).\footnote{La formule est publiée par Cardan
dans l'\emph{Ars magna} (1545), qui la tient de del Ferro et de Tartaglia ; elle
porte aujourd'hui son nom. Le traité ne cite personne ici.}

Le traité exige, au second cas, que $D^2$ surpasse $4B^3$. C'est la bonne
condition : $u^3$ et $v^3$ sont les racines de $t^2 - Dt + B^3 = 0$, qui sont
réelles et positives si et seulement si $D^2 \ge 4B^3$. L'égalité, que le
traité n'envisage pas, donne $u = v$ et une racine double $x = 2u$.\footnote{En
tête du chapitre, la même condition est dite en mots : « quadruplus cubus è
triente coefficientis plani \emph{cedit} solidi dati \emph{medietatis}
quadrato » (vue 13), et de même à la vue 14 avec « præstat » (italiques : lectures
incertaines de la transcription). Telle que la phrase est lue, elle compare
quatre fois $(p/3)^3$ au carré de la \emph{moitié} du solide, $(q/2)^2$, où $p$ est
le coefficient et $q$ le second membre ; la bonne condition est soit
$(q/2)^2 \gtrless (p/3)^3$, soit $q^2 \gtrless 4(p/3)^3$, et c'est la seconde que
donne l'énoncé du théorème II (« præstet autem $D$ solidi quadratum quadruplo $B$
plani cubo »). Le mot « medietatis » est donc de trop, ou mal lu.}

\subsection*{La condition démontrée par une seconde main}

Une main rapide, à l'encre noire, ajoute au pied de la vue 13 une preuve de
cette condition, avec en marge « vide Zet. 20 lib. 2. » : si les côtés sont
$u$ et $v$, alors $(u^3 + v^3)^2 = u^6 + 2u^3v^3 + v^6 > 4u^3v^3$, parce que
$u^6$, $u^3v^3$, $v^6$ sont en proportion continue et que, de trois
proportionnelles inégales, la somme des extrêmes surpasse le double de la
moyenne.\footnote{C'est l'inégalité entre moyenne arithmétique et moyenne
géométrique, $a + c > 2\sqrt{ac}$ pour $a \ne c$ positifs. Le renvoi « vide Zet. 20
lib. 2. » vise un vingtième zététique d'un deuxième livre que ce volume ne
contient pas ; il n'a pas été suivi.}

La même main démontre, sur un béquet marqué « 5 bis » (vue 15, « vel sic
ἀποδεικτικῶς », apodeiktikôs, « démonstrativement ») et dans une première rédaction barrée en marge de la vue 16, la
condition du troisième théorème. Pour deux côtés $A$ et $B$, on pose
$p = A^2 + AB + B^2$ et $q = AB(A + B)$, et il s'agit de montrer que
$4(p/3)^3 > q^2$. La main développe, retranche les termes communs et trouve
qu'il faut
\[ 4A^6 + 12A^5B + 12AB^5 + 4B^6 \;>\; 3A^4B^2 + 26A^3B^3 + 3A^2B^4 , \]
ce qu'elle établit en comparant des termes de la suite continue
\[ A^6,\; A^5B,\; A^4B^2,\; A^3B^3,\; A^2B^4,\; AB^5,\; B^6 . \] Le développement est juste, et la
preuve aussi : la différence des deux membres est
\[ 4p^3 - 27q^2 = (A - B)^2\,(2A + B)^2\,(A + 2B)^2 , \]
qui est, au signe près, le discriminant de la cubique $x^3 - px - q$ — positif
dès que $A \ne B$.\footnote{Le calcul est de cette lecture ; la main ne factorise
pas, elle compare. Le mot grec de la vue 15 est lu avec doute ; après
« hypothesi », une tache couvre un ou deux signes. Le béquet porte en tête un
signe de renvoi, une croix surmontée d'un petit cercle suivie de « 4 », qui se
retrouve dans la ligne et dans la marge de la vue 16.}

\subsection*{Le cas irréductible}

Le troisième théorème traite le cas où $x^3 - 3Bx = D$ avec $D^2 < 4B^3$ :
alors l'équation « inverse » $3By - y^3 = D$ a aussi des racines, et
\begin{itemize}
  \item il y a deux côtés $u$ et $v$ tels que $u^2 + uv + v^2$ vaut le
    coefficient $3B$ et que $uv(u + v) = D$ ;
  \item $x = u + v$ est la racine de $x^3 - 3Bx = D$, et $u$, $v$ sont les deux
    racines de $3By - y^3 = D$.
\end{itemize}
C'est juste : $(u+v)^3 - (u^2+uv+v^2)(u+v) = uv(u+v)$, et de même pour $u$ et
$v$ dans l'équation inverse. L'exemple l'est aussi : $x^3 - 21x = 20$ donne
$5$, et $21y - y^3 = 20$ donne $1$ ou $4$, avec $1 + 4 + 16 = 21$ et
$1 \cdot 4 \cdot 5 = 20$.\footnote{Le feuillet dit que les carrés des deux côtés,
ajoutés à leur rectangle, « font $B$ planum » ; il faut « $B$ planum ter », le
coefficient lui-même, comme le montre l'exemple du feuillet, où la somme vaut
$21$ et non $7$. Dans l'exemple, « $21N - 1Q$ » est corrigé par l'écrivain en
« $21N - 1C$ ». Le paragraphe suivant (« Neq; vero in Geometrica phrasi\ldots ») dit
la même chose en proportions : $1 + 4 + 16 = 21$ est la somme des carrés de
trois proportionnelles $1, 2, 4$, et $20 = (1 + 4) \cdot 2^2$. En termes
d'aujourd'hui, l'équation $x^3 - 21x - 20 = 0$ a trois racines réelles, $5$,
$-1$ et $-4$ ; celles de l'équation inverse en sont les opposées.}

Le traité juge alors plus élégante une constitution tirée « des sections
angulaires », et la donne sous le titre « Aliter, tertium theorema ». Si
$x^3 - 3B^2x = B^2D$ avec $B > D/2$, on prend deux triangles rectangles de même
hypoténuse $B$, l'angle aigu du premier triple de celui du second, et tels que
le double de la base du premier soit $D$ ; alors $x$ est le double de la base
du second, et les racines de l'équation inverse $3B^2y - y^3 = B^2D$ sont la base
du second plus ou moins $\sqrt{3}$ fois sa hauteur. C'est la résolution
trigonométrique de la cubique : si $x = 2B\cos\theta$,
\[ x^3 - 3B^2x = 2B^3\cos 3\theta , \]
de sorte que l'équation revient à $\cos 3\theta = D/(2B)$, qui a une solution
si et seulement si $D < 2B$ ; les deux autres racines sont
$2B\cos(\theta \pm 120^\circ) = -B\cos\theta \mp \sqrt{3}\,B\sin\theta$, dont les
opposées sont les racines positives de l'équation inverse.\footnote{C'est
le cas que l'on appelle aujourd'hui \emph{casus irreducibilis} : trois racines
réelles, et une formule de Cardan qui passe par des racines carrées de nombres
négatifs. Le nom est postérieur au traité ; la condition $D^2 < 4B^3$ du
théorème III, écrite pour $x^3 - 3Bx = D$, et la condition $B > D/2$ de
l'« Aliter », écrite pour $x^3 - 3B^2x = B^2D$, sont la même.} L'exemple est juste :
pour $x^3 - 300x = 432$, l'hypoténuse est $10$ ; le second triangle a pour base
$9$ et pour hauteur $\sqrt{19}$ ; le premier a pour base $10 \cos 3\theta
= 2{,}16$, avec $\cos\theta = 0{,}9$ ; $x = 18$, et $300y - y^3 = 432$ donne
$9 \pm \sqrt{57}$.\footnote{$\cos 3\theta = 4 \cdot 0{,}729 - 3 \cdot 0{,}9 =
0{,}216$ ; $18^3 - 300 \cdot 18 = 432$ ; $\sqrt{3 \cdot 19} = \sqrt{57}$. La page
écrit « basis autem primi dupla est $\frac{432}{100}$ » et « fit basis
$2\frac{16}{100}$ ». Dans « $1C - 300N$ », le $N$ est ajouté au-dessus de la
ligne. Le mot entre « contracta » et « longitudine », dans l'énoncé, n'est pas
lu.}

La seconde main démontre cette constitution dans les marges des vues 17 et 18,
sur la figure d'un cercle de diamètre $AB$ dont l'arc $FB$ est partagé en trois
en $G$ et $K$. En posant $AB = B$, $2\,AG = A$ et $2\,AF = D$, elle obtient
$A^3 - 3B^2A = B^2D$ par deux proportions successives ; puis, prolongeant $BG$
jusqu'à $L$ de sorte que $BL = 2\,BG$, elle traite l'équation inverse. C'est la
relation $\cos 3\alpha = 4\cos^3\alpha - 3\cos\alpha$ écrite sur les cordes issues
de l'extrémité d'un diamètre ($AG = B\cos\alpha$, $AF = B\cos 3\alpha$), et le
calcul qu'on peut suivre est juste.\footnote{La lecture de ces deux longues
marges est donnée par la transcription « sous réserve » : les fractions y sont
serrées, plusieurs lettres de la figure sont reprises, un passage s'arrête au
pied de la vue 17 sur « eiusdem » et reprend à la vue 18. Cette lecture n'en
restitue pas chaque pas.}

Cette main parle d'elle-même : « quoniam demonstratum a nobis ad sectiones
angulares », puis « est è demonstratis a nobis \emph{in tractatu de sectionibus
angularibus} ».\footnote{« Puisqu'il a été démontré par nous aux sections
angulaires », « il suit de ce que nous avons démontré dans le traité des
sections angulaires » ; les mots en italique sont une addition interlinéaire lue
avec doute. Un traité des sections angulaires est connu dans la littérature
sous le nom de Viète ; la coïncidence du sujet ne suffit pas à identifier cette
main, et cette lecture ne l'identifie pas.} Le texte principal, de son côté,
renvoie ce dernier théorème à « un recueil particulier de zététiques qui
touchent aux sections angulaires », et clôt le chapitre.\footnote{« hoc postremo
theoremate excepto, quod par est reiici in peculiarem Zeteticorum ad angulares
sectiones pertinentium \ldots{} » (vue 17) ; le dernier mot, souligné, n'est pas
lu, ni le mot après « Expositis ».}

\section*{Changer d'inconnue : les chapitres VII et VIII}

\pagerange{18}{23}
Le chapitre VII pose la doctrine sur laquelle repose tout le reste des deux
traités : transformer une équation. On peut le faire en changeant l'inconnue,
ou en la gardant.

\subsection*{Changer l'inconnue}

La nouvelle inconnue $y$ doit avoir avec l'ancienne « une différence ou une
raison donnée », de sorte que connaître l'une donne l'autre. Le traité énumère
sept manières :

\begin{itemize}
  \item première, par addition, $y = x + B$ ;
  \item deuxième, par soustraction, $y = x - B$ ou $y = B - x$ ;
  \item troisième, par multiplication, $y = Bx$ ;
  \item quatrième, par division, $y = x/B$ ;
  \item cinquième et sixième, par proportion, $B : G = x : y$ ou
    $x : B = G : y$, c'est-à-dire $x = By/G$ ou $x = BG/y$ ;
  \item septième, par une relation du second degré entre les deux inconnues
    (la « parabolica hypostasis ») : $y^2 + xy = D$, $y^2 - xy = D$ ou
    $xy - y^2 = D$, d'où $x = (D - y^2)/y$, $x = (y^2 - D)/y$ ou
    $x = (y^2 + D)/y$ ;
\end{itemize}

et enfin par des manières composées, « à imaginer et à essayer ».
Dans chaque cas on remplace $x$ par sa valeur en $y$ et l'on « ordonne »
l'équation nouvelle. Les six premières manières sont les substitutions affines
et l'inversion ; la septième est une substitution rationnelle, et c'est
exactement celle dont le second traité se servira pour résoudre la cubique
(vues 119 à 122) : $x = B/y - y$, qu'on appelle aujourd'hui la substitution de
Viète.\footnote{Le nom est moderne, et il est donné ici parce que l'énoncé
coïncide : pour $x^3 + 3Bx = 2Z$, la substitution $x = (B - y^2)/y$ ramène
l'équation à une équation du second degré en $y^3$. Tschirnhaus (1683)
généralisera l'idée de changer d'inconnue pour faire disparaître des termes, en
prenant pour nouvelle inconnue un polynôme en l'ancienne ; les substitutions de
ces chapitres n'en sont que le cas le plus simple, et le traité ne va pas
au-delà.}

\subsection*{Garder l'inconnue}

Sans changer l'inconnue, on peut élever ou abaisser le degré de l'équation : la
« montée climactique » régulière élève les deux membres au carré ou au cube ;
l'irrégulière multiplie tous les termes par une même puissance de l'inconnue,
éventuellement affectée ; la descente fait l'inverse, par division, « en
ajoutant le supplément donné ou à chercher ». La justification tient en une
phrase : des grandeurs égales ont des puissances égales, et un multiplicateur
ou un diviseur commun ne change pas l'égalité.\footnote{« quoniam quæ æqualia
sunt longitudine, æqualia quoq; sunt similis potestate, et contra. neq;
communis divisor, vel multiplicator æqualitatem Immutat vel rationem »
(vue 20). La réciproque (« et contra ») n'est vraie que pour des grandeurs
positives, ce que le traité suppose toujours. La vue 19 porte un mot non lu
après « in » (« condicibus » ou « ordinibus »).}

\subsection*{Le plasme}

Le chapitre VIII appelle \emph{plasme} la génération d'une équation affectée à
partir d'une équation plus pure, par l'un de ces changements, et pose quatre
règles.\footnote{Le chapitre commence à la vue 20 et se poursuit à la vue 21 ;
la phrase de la vue 20 s'arrête sans ponctuation sur « si quo id liceat
compendio », et la vue 21 commence par « Quod cum deprehendit \emph{sic
mutatarius} », lecture que la transcription donne pour douteuse. Une petite main
pâle a posé dans la marge de la vue 20 deux notes appelées par des croix,
« + in exemplari \emph{excripto} / vel potius \emph{adscriptum} » et « + vel
in », lues en partie.}

\begin{enumerate}
  \item Une équation dont tous les degrés intermédiaires sont présents provient,
    par addition ou soustraction, d'une équation où manque le terme qui suit la
    plus haute puissance : la racine a été augmentée ou diminuée de la moitié du
    coefficient de $x$ dans le second degré, du tiers du coefficient de $x^2$
    dans le troisième, du quart de celui de $x^3$ dans le quatrième, du
    cinquième dans le cinquième, « et ainsi de suite ». C'est la réduction à
    l'\emph{équation privée de son second terme} : la substitution
    $x = y - a_{n-1}/n$ fait disparaître le terme en $y^{n-1}$ de
    $x^n + a_{n-1}x^{n-1} + \cdots$. Ainsi les quadratiques affectées viennent
    des pures, les cubiques affectées sous le carré et le côté viennent des
    cubiques affectées sous le seul côté, et de même plus haut.
  \item Une équation dont la racine est « plane » ou « solide » (où l'inconnue
    n'entre que par son carré ou son cube) provient d'une autre par
    multiplication ou par proportion.\footnote{Au-dessus de « homogenea »,
    l'écrivain propose « climactica potius ». La marge de la vue 21 en donne
    deux exemples : si $x^2 + Bx = D$ et $Bx = y^2$, alors $y^4 + B^2y^2 = B^2D$, et
    $y$ est moyenne proportionnelle entre $B$ et $x$ ; si $x^2 + Bx = D^2$ et
    $B^2x = y^3$, alors $y^6 + B^3y^3 = B^4D^2$. Les deux calculs sont justes. Le
    premier est écrit avec « $D$ \emph{solido} », que la dimension interdit (il
    faut un plan), et sa conclusion se lit « \emph{Erit} itaque E. Cubi affecti
    quadratoquadrati media proportionalis » là où l'on attend « $E$ latus » : ce
    sont des lectures de lettres.}
  \item Toute équation du quatrième degré dont manquent certains termes provient
    d'une équation du second degré par « montée climactique » : on répartit les
    termes de la quadratique entre les deux membres et l'on élève au carré.
  \item Toutes les affections cubiques se déduisent des quadratiques par montée
    régulière, mais on ne peut en revenir que par des équations de même degré.
\end{enumerate}

\subsection*{Les « anomales plasmatiques », et un porisme}

Le chapitre remarque enfin (vue 23) que si l'on forme $y = B - x$ avec $B$ plus
grand que la racine, à partir de $x^2 = Z$, on tombe toujours sur une équation
de la forme $2By - y^2 = B^2 - Z$ — celle dont les deux racines $B \pm \sqrt{Z}$
sont positives. D'où le porisme : « toutes les fois que la puissance est
retranchée de l'homogène d'affection, la racine est ambiguë ». L'énoncé, pris
tel quel, est trop fort. Ce qui est vrai : une équation de la forme
$(\text{termes affectés}) - x^n = c$, avec $c > 0$, a zéro ou deux racines
positives (comptées avec leur multiplicité) quand ses termes affectés sont d'un
seul signe ; elle en a deux si elle en a une.\footnote{Par la règle des signes de
Descartes (1637), postérieure au traité : la suite des coefficients de
$x^n - (\ldots) + c$ présente alors exactement deux changements de signe, et le
nombre de racines positives est $2$ ou $0$. Ainsi $Bx - x^3 = c$ a deux racines
positives si $c < 2(B/3)^{3/2}$ et aucune au-delà. Le chapitre XV du traité
(vues 45 et 46) démontre l'ambiguïté dans le cas quadratique, où elle est
exacte.}

\section*{Les plasmes : chapitres IX à XIII}

\pagerange{23}{43}
Ces cinq chapitres sont une collection de théorèmes, chacun démontré en
développant et en « ordonnant », la plupart suivis d'un exemple numérique. Ils
appliquent systématiquement les changements d'inconnue du chapitre VII. On les
donne ici dans l'ordre, en langage d'aujourd'hui ; tous ont été refaits, et ils
sont justes, sauf là où une note dit le contraire.

\subsection*{Chapitre IX : les quadratiques tirées des pures}

Si $x^2 = Z$,
\begin{itemize}
  \item $y = x + B$ donne $y^2 - 2By = Z - B^2$ ;
  \item $y = x - B$ donne $y^2 + 2By = Z - B^2$ ;
  \item $y = B - x$, ou $y = x + B$, donne $2By - y^2 = B^2 - Z$.
\end{itemize}
Au premier théorème, le second membre est écrit avec le double trait,
« $Z$ plano $=$ $B$ quadrato » : $Z$ et $B^2$ peuvent être dans un ordre ou dans
l'autre, et c'est ce que dit la différence signée $Z - B^2$.\footnote{C'est le
premier emploi, dans ce volume, du double trait « = » pour une différence
d'ordre non donné ; voir la section sur la notation. La transcription note que
le théorème II, à la même place, écrit un moins simple. Sur le feuillet, le numéro
du chapitre était d'abord « VIII », biffé et récrit « IX ».}

\subsection*{Chapitre X : cubiques affectées sous le carré, tirées des cubiques affectées sous le côté}

Si $x^3 - 3B^2x = Z$, la translation fait apparaître un terme en $y^2$ et fait
disparaître celui en $y$ :
\[ y = x - B : \; y^3 + 3By^2 = Z + 2B^3 ; \qquad
   y = x + B : \; y^3 - 3By^2 = Z - 2B^3 . \]
Si $3B^2x - x^3 = Z$ :
\[ y = B - x : \; 3By^2 - y^3 = 2B^3 - Z ; \qquad
   y = B + x : \; 3By^2 - y^3 = 2B^3 + Z . \]

\subsection*{Chapitre XI : cubiques affectées sous le carré et le côté}

Neuf théorèmes traitent $x^3 + Dx = Z$, $x^3 - Dx = Z$ et $Dx - x^3 = Z$, chacun
par $y = x + B$, $y = x - B$ et $y = B - x$. Par exemple, $x^3 + Dx = Z$ et
$y = x + B$ donnent
\[ y^3 - 3By^2 + (3B^2 + D)\,y = Z + B^3 + DB , \]
et $Dx - x^3 = Z$ avec $y = B - x$ donne
\[ -y^3 + 3By^2 + (D - 3B^2)\,y = DB - B^3 - Z . \]
Les neuf développements ont été refaits ; tous sont justes, une fois les signes
corrigés par l'écrivain lui-même.\footnote{Aux vues 28 et 32, des « $+$ »
barrés d'une croix deviennent des « $-$ » ; à la vue 30, le second membre du
quatrième théorème porte le double trait devant $DB$ ; à la vue 32, le cube à
retrancher est écrit avec deux colonnes de signes, dont l'intérieure est la
bonne, et il manque un « $-$ » devant « $E$ cubo ». La marge de la vue 32 porte,
d'une écriture plus petite et plus pâle, des exemples chiffrés :
« $19N - 1C$ æq. $30$ », dont les racines positives sont $2$ et $3$ ;
« $52N - 1C$ æq. $96$. fit N $2$ vel $6$ », qui est juste (et $2^2 + 2 \cdot 6 + 6^2
= 52$, $2 \cdot 6 \cdot 8 = 96$, comme au chapitre VI) ; et
« $1C + 9Q - 24N$ æq. $33$ », qui, telle qu'elle est lue, n'a pas de racine
entière (sa racine positive est voisine de $2{,}946$) : une lecture ou une
écriture fautive, que cette lecture ne restitue pas. Une petite figure au pied
de la marge — un demi-cercle sur un diamètre partagé en $2$ et $6$, la
perpendiculaire cotée « $\sqrt{12}$ » et la corde « $\sqrt{16}$ » — est la
construction de la moyenne proportionnelle $\sqrt{2 \cdot 6}$. Le premier mot de
la marge, en lettres grecques, n'est pas lu.}

\subsection*{Chapitre XII : racines planes et solides}

Si $x^2 + Bx = Z$ et $Bx = y^2$, alors $y^4 + B^2y^2 = B^2Z$ ; si $B^2x = y^3$,
alors $y^6 + B^3y^3 = B^4Z$ ; si $x^3 + Bx^2 = Z$ et $Bx = y^2$, alors
$y^6 + B^2y^4 = B^3Z$ ; et de même pour les deux autres espèces de chaque
équation. La nouvelle inconnue est une moyenne proportionnelle entre $B$ et
$x$, ou la seconde de quatre proportionnelles de premier terme $B$ et de
quatrième $x$.\footnote{Deux corrections de l'écrivain, toutes deux justes : à la
vue 34, « $B$ in $D$ » corrigé en marge en « $B$ quadrato » ; à la vue 35,
« omnia ducantur in $B$ quadratum » corrigé en « quad. quad. ». La conclusion
du théorème IIII écrit « $E$ solidum » là où l'énoncé a « $E$ cubum » : c'est
le même objet, nommé par sa dimension.}

\subsection*{Chapitre XIII : quartiques tirées des quadratiques}

C'est ici qu'apparaît l'idée qui fera la force du second traité. Si
$x^2 + Bx = Z$, on multiplie, on substitue $x^2 = Z - Bx$, et l'on obtient une
équation du quatrième degré qui a la même racine :
\[ x^4 + (B^3 + 2BZ)\,x = Z^2 + B^2Z . \]
Exemple juste : $x^2 + 8x = 20$ donne $x^4 + 832x = 1680$, de racine $2$. De même
pour $x^2 - Bx = Z$ (racine $10$ de $x^2 - 8x = 20$, et $x^4 - 832x = 1680$) et
pour $Bx - x^2 = Z$ ($12x - x^2 = 20$, $1248x - x^4 = 2480$, racine $2$).
Le théorème II, sauté sur la page, est sur une fiche montée (vue 38), qui porte
« ad faciem aversam foly 16. » et un astérisque répondant à celui de la
vue 37.\footnote{La quartique n'a pas que la racine de la quadratique dont elle
vient : $x^4 + (B^3 + 2BZ)x - Z^2 - B^2Z = (x^2 + Bx - Z)(x^2 - Bx + Z + B^2)$, et
le second facteur n'a pas de racine réelle. Le traité ne le dit pas, mais il
n'en a pas besoin : il ne prétend pas que la quartique n'ait que cette racine.}

Les théorèmes IIII à XII font de même pour la quadratique dont le second membre
est une somme ou une différence, $x^2 + Bx = S \pm D$, en portant $\mp D$ du côté de
$x^2$ et en élevant au carré :
\[ (x^2 - D)^2 = (S - Bx)^2 \;\Longrightarrow\;
   x^4 - (2D + B^2)\,x^2 + 2BSx = S^2 - D^2 , \]
puis toutes les combinaisons de signes. Les exemples sont justes : pour
$x^2 + 8x = 20 = 15 + 5$, $x^4 - 74x^2 + 240x = 200$ ; pour $12x - x^2 = 20$,
$114x^2 - 120x - x^4 = 200$ ; pour $x^2 + 8x = 20 = 60 - 40$,
$x^4 + 16x^2 + 960x = 2000$ ; pour $x^2 - 8x = 20 = 60 - 40$,
$x^4 + 16x^2 - 960x = 2000$ (racine $10$) ; pour $12x - x^2 = 20 = 60 - 40$,
$960x + 24x^2 - x^4 = 2000$.\footnote{Au théorème V, le renvoi « in antecedente
capite » vise le théorème précédent du même chapitre. Au théorème IX, la
transcription hésite sur les valeurs de $S$ et $D$ ; l'exemple est cohérent avec
$S = 60$ et $D = 40$ partout ($2 \cdot 12 \cdot 40 = 960$, $144 - 120 = 24$,
$3600 - 1600 = 2000$). Au même théorème, « antecedente » est biffé et remplacé
par « septimo ».}

Les théorèmes X à XII (« De quadrato quadratis adfectis sub cubo », vues 41 à
43) éliminent au contraire le terme en $x^2$. De $x^2 + Bx = Z$ on tire
$x^3 = (Z + B^2)x - BZ$, donc $x = (x^3 + BZ)/(Z + B^2)$ et
$x^2 = (Z^2 - Bx^3)/(Z + B^2)$ ; en portant cette valeur dans le carré de
l'équation, on obtient
\[ x^4 + \frac{2BZ + B^3}{Z + B^2}\,x^3 = \frac{Z^3}{Z + B^2} . \]
L'exemple est juste : $x^2 + 14x = 147$ donne $x^4 + 20x^3 = 9261$, de racine
$7$, car $(2 \cdot 14 \cdot 147 + 14^3)/343 = 20$ et $147^3/343 = 9261$. De même,
$x^2 - 14x = 147$ donne $x^4 - 20x^3 = 9261$, de racine $21$ ; et, pour
$Bx - x^2 = Z$,
\[ \frac{B^3 - 2BZ}{B^2 - Z}\,x^3 - x^4 = \frac{Z^3}{B^2 - Z} , \]
avec l'exemple $21x - x^2 = 98$, $15x^3 - x^4 = 2744$, de racine $7$.\footnote{Aux
théorèmes X et XI, le trait de division est tiré sous tout le premier membre,
$x^4$ compris ; il ne doit porter que sur le coefficient de $x^3$, comme le
montrent les exemples et comme l'écrit la note marginale du théorème XII. À
l'énoncé du théorème X (vue 41), le second membre est écrit « $Z$ plano
plano » sur « $Z$ plano $+ B$ quadrato » ; la fin de la démonstration (vue 42)
porte « $Z$ plano plano plano », qui est juste ($Z^3$ pour un plan $Z$). La note
marginale de la vue 43, « ad sinistram paginam fol. 18. », justifie la valeur
de $x^2$ en renvoyant au feuillet 18 de la série à l'encre (vue 41). Le premier
chiffre de « 2744 » est écrit sur un 7. La marge de la vue 42 porte des calculs
et quatre petites figures (demi-cercles, triangle rectangle) qui ne sont pas
interprétés ici.}

\section*{Des quadratiques aux cubiques ; l'ambiguïté démontrée}

\pagerange{43}{46}
Le chapitre XIIII multiplie une quadratique par $x$ et remplace $x^2$ par sa
valeur :
\[ x^2 + Bx = Z \;\Longrightarrow\; (Z + B^2)\,x - x^3 = BZ , \]
avec l'exemple $x^2 + 8x = 20$, $84x - x^3 = 160$, de racine $2$. Les théorèmes
suivants donnent $x^3 - (B^2 + Z)x = BZ$ pour $x^2 - Bx = Z$ (racine $10$),
$(B^2 - Z)x - x^3 = BZ$ pour $Bx - x^2 = Z$ ($124x - x^3 = 240$, racine $2$),
puis, pour un second membre $BD$, $x^3 + (B + D)x^2 = BD^2$ ($x^2 + 16x = 80$,
$x^3 + 21x^2 = 400$, racine $4$), $(B + D)x^2 - x^3 = BD^2$ (racine $20$) et
$(B - D)x^2 - x^3 = BD^2$ ($9x - x^2 = 18$, $7x^2 - x^3 = 36$, racine $3$). Tous
sont justes.\footnote{La cubique déduite a d'autres racines que celle de la
quadratique : au premier théorème,
$x^3 - (Z + B^2)x + BZ = (x - B)(x^2 + Bx - Z)$, de sorte que $84x - x^3 = 160$ a
aussi la racine positive $8 = B$ ($672 - 512 = 160$). Le traité ne le dit pas ;
c'est un ajout. À la vue 44, « $1C - 84$ » est écrit sans le signe N.}

Le chapitre XV démontre, dans le cas quadratique, ce que le porisme de la
vue 23 affirmait en général : si $x$ diffère de $B$ de $S$, dans un sens ou dans
l'autre, alors dans les deux cas
\[ 2Bx - x^2 = B^2 - S^2 , \]
et la racine est « double », $B - S$ ou $B + S$. Exemple : $B = 6$, $S = 4$,
$12x - x^2 = 20$, de racines $2$ et $10$. Le traité conclut que toute équation
$Dx - x^2 = Z$ se lit ainsi, avec $D = 2B$ et $Z = B^2 - S^2$.\footnote{La page écrit
« $Z$ planum esse $S$ quadratum minus $B$ quadrato » : l'ordre des deux carrés est
inversé ; l'égalité démontrée juste au-dessus donne $B^2 - S^2$. La fin de la phrase
se perd dans le pli. Le numéro du chapitre, d'abord « XIIII », est biffé et récrit « XV ». Le paragraphe
qui suit renvoie, pour les cubes et les carrés-carrés « négatifs », au sixième
théorème du chapitre précédent et aux théorèmes III, VI et IX d'un chapitre écrit
« noni », que la marge corrige en « duodecimi » : ce sont ceux du chapitre XIII
de la numérotation actuelle, qui était le XII\textsuperscript{e} avant que
l'écrivain ne décale ses numéros d'une unité (vues 23 et 25). C'est une
inférence de cette lecture.}

\section*{La syncrise : les coefficients par les racines}

\pagerange{46}{55}
Le chapitre XVI définit la \emph{syncrise} : la comparaison de deux équations
« corrélatives », c'est-à-dire semblables et faites des mêmes données, mais qui
ont des racines différentes, soit parce que la forme de l'équation admet deux
racines, soit parce que les signes diffèrent. Il en distingue trois espèces
pour les équations à un seul terme affecté, écrites ici avec un exposant
$n$ pour la puissance et un exposant $k < n$ pour le degré affecté :

\begin{itemize}
  \item les \emph{ambiguës}, où la puissance est retranchée dans les deux :
    $Bx^k - x^n = Z$ et $By^k - y^n = Z$ ;
  \item les \emph{contradictoires}, où elle est ajoutée dans l'une et retranchée
    dans l'autre : $x^n + Bx^k = Z$ et $y^n - By^k = Z$ ;
  \item les \emph{inverses}, où la puissance est diminuée du terme affecté dans
    l'une et le terme affecté diminué de la puissance dans l'autre :
    $x^n - Bx^k = Z$ et $By^k - y^n = Z$.\footnote{Le nom de la troisième espèce
    est lu « inversarum » d'après les lettres (« Jnuersarum »), aux vues 46 et
    47, et « Inversis » ensuite ; la transcription le donne pour douteux aux deux
    premières occurrences.}
\end{itemize}

La méthode est toujours la même : deux choses égales à une même troisième sont
égales entre elles ; on égale les deux premiers membres, on fait passer les
puissances d'un côté et les termes affectés de l'autre, et l'on divise par la
différence ou la somme des degrés affectés. On obtient ainsi le coefficient et
le terme connu en fonction des deux racines. Pour les ambiguës, avec $x > y$ :
\[ B = \frac{x^n - y^n}{x^k - y^k}, \qquad
   Z = \frac{x^n y^k - y^n x^k}{x^k - y^k} ; \]
pour les contradictoires,
\[ B = \frac{y^n - x^n}{y^k + x^k}, \qquad
   Z = \frac{x^n y^k + y^n x^k}{y^k + x^k} ; \]
pour les inverses,
\[ B = \frac{x^n + y^n}{x^k + y^k}, \qquad
   Z = \frac{x^n y^k - y^n x^k}{x^k + y^k} ; \]
et pour les inverses dont l'une est affectée par addition,
$x^n + Bx^k = Z$ et $By^k - y^n = Z$,
\[ B = \frac{x^n + y^n}{y^k - x^k}, \qquad
   Z = \frac{x^n y^k + y^n x^k}{y^k - x^k} . \]
Les quatre paires de formules ont été vérifiées ; elles sont justes.\footnote{Le
feuillet écrit plusieurs signes fautifs, que ses propres marges corrigent pour
la plupart. Au problème II (vue 51), le diviseur porte deux fois le double trait
« $=$ » là où il faut « $+$ », que la marge rétablit. Au problème III (vues 52 et
53), l'énoncé porte le double trait devant le terme affecté, là où l'antithèse
qui suit suppose la soustraction ; au diviseur, la page écrit « $A$ gradu $+ B$
gradu », avec en marge « $- E.$ » : il faut « $+ E$ gradu », comme l'écrit la
ligne précédente. Au dernier cas (vues 53 et 54), « $A$ potestatem $- E$ potestate »
est corrigé en marge par « $+$ », et deux signes sont laissés en blanc (un point),
avec « $+$ » en marge. Les mots lus « Iucunda » (vue 47), « Divinus »,
« lucidum », « idq; », « transitus », « retentis », « pugnas » (vue 54) sont
donnés pour douteux par la transcription, et plusieurs phrases de ces pages
n'en sont pas sûres ; la substance mathématique ne l'est pas moins, car chaque
formule est écrite en espèces.}

Dans le cas le plus simple, $Bx - x^2 = Z$ et $By - y^2 = Z$, le traité trouve
\[ B = x + y, \qquad Z = xy , \]
et, pour $Bx - x^3 = Z$ et $By - y^3 = Z$,
\[ B = x^2 + xy + y^2, \qquad Z = xy\,(x + y) . \]
Ce sont, pour ces deux formes, les \emph{relations entre coefficients et racines}
que l'on appelle aujourd'hui les formules de Viète : le coefficient est la somme
des racines, le terme connu leur produit.\footnote{Le nom est moderne, et
l'énoncé coïncide ici pour les racines positives que le traité considère. Pour
la cubique, $x^3 - Bx + Z = 0$ a une troisième racine, $-(x + y)$, négative, que
le traité ne nomme pas ; avec elle, la somme des trois racines est nulle, la
somme de leurs produits deux à deux vaut $-B$ et leur produit $-Z$, comme le
veulent les relations générales. La forme la plus générale est au chapitre
XIIII du second traité (vues 139 à 141). Au cas cubique, la page écrit « $D$
plano » là où le problème porte « $B$ plano », et la marge corrige par des
« $B$ » ; elle écrit aussi « omnibus per $A$ divisis » là où le diviseur écrit
dessous est $A - E$.} Le traité en tire aussitôt des « formules » d'équations
dont on choisit les racines : pour les racines $F$ et $G$,
\[ (F + G)\,x - x^2 = FG , \qquad (F^2 + FG + G^2)\,x - x^3 = FG(F + G) , \]
avec $F = 10$, $G = 2$ : $12x - x^2 = 20$ et $124x - x^3 = 240$, qui ont bien
pour racines $10$ et $2$.

La fin du chapitre dit pour quelles puissances chaque espèce est
« efficace », c'est-à-dire donne des coefficients entiers en les racines, et la
raison qu'il en donne est exacte : $x^n - y^n$ se divise toujours par $x - y$ ;
$x^n - y^n$ se divise par $x + y$ quand $n$ est pair ; $x^n + y^n$ se divise
par $x + y$ quand $n$ est impair ; et $x^n + y^n$ ne se divise jamais par
$x - y$. Les ambiguës sont donc efficaces à tous les degrés, les contradictoires
aux degrés pairs, les inverses aux degrés impairs.\footnote{Le traité renvoie ces
faits à un « Caput de genesi potestatum a binomia radice » — un chapitre sur la
génération des puissances d'un binôme — que ce volume ne contient pas. La
marge de la vue 54 écrit les quotients pour $n = 2$, $3$, $4$ : $x + y$,
$x^2 + xy + y^2$, $x^3 + x^2y + xy^2 + y^3$, et renvoie, « pro cubis », à un
chapitre dont le numéro n'est lu qu'en partie (« Cap. \ldots{} X. lib. 2. »). La
première phrase de la vue 55 (« Ceterum dum aggregatum potestatum applicatur ad
differentiam radicum nulla inde orietur magnitudo\ldots ») énonce le dernier des
quatre faits.}

\section*{La « phrase géométrique » : suites de proportionnelles}

\pagerange{55}{58}
Le chapitre XVII habille la syncrise en géométrie, par sept propositions sur les
suites de trois à sept proportionnelles $a_1, a_2, \ldots, a_n$ de raison $r$.
Toutes sont justes, parce que chaque terme vaut le précédent multiplié par
$r$.\footnote{La marge de la vue 57 ajoute que $BG^2 - B^2G = F^3 - D^3$, ce qui est
juste. La marge de la vue 55 porte un calcul très corrigé, avec des « plans
proportionnels » $F^2$, $FE$, $E^2$, que la transcription ne reconstitue pas ; il
ne l'est pas ici.}
\begin{itemize}
  \item $a_1 : a_n = a_1^{\,n-1} : a_2^{\,n-1}$ (proposition I) ;
  \item $a_1 : a_n$ est aussi le rapport de la somme des puissances
    $(n-1)$-ièmes des $n - 1$ premiers termes à celle des $n - 1$ derniers
    (II), le rapport des puissances $(n-1)$-ièmes de leurs sommes (III), celui
    des sommes alternées de leurs puissances (IIII), celui des puissances de
    leurs sommes alternées (V) ;
  \item pour quatre proportionnelles $B, D, F, G$ (VI et VII) :
    \[ G^3 + 3D^3 - B^3 - 3F^3 = (G - B)^3 , \qquad
       G^3 + 3D^3 + B^3 + 3F^3 = (G + B)^3 , \]
    parce que $D^3 = B^2G$ et $F^3 = BG^2$.
\end{itemize}

\section*{Les constitutions : ambiguës, contradictoires, inverses}

\pagerange{58}{68}
Les chapitres XVIII à XXI donnent, pour chaque espèce, la constitution des
équations à partir des termes d'une suite de proportionnelles, avec pour chaque
proposition un exemple dont les racines sont connues d'avance. C'est la
syncrise appliquée : on choisit la suite, on en tire les coefficients.

\subsection*{Chapitre XVIII : les ambiguës}

Pour les ambiguës « affectées sous le côté », $Bx - x^n = Z$ :
\begin{itemize}
  \item $n = 2$ : $B$ est la somme de deux côtés, $Z$ leur produit ; exemple
    $12x - x^2 = 20$, racines $2$ et $10$ ;
  \item $n = 3$ : $B$ est la somme des carrés de trois proportionnelles
    $a, b, c$, et $Z = a(b^2 + c^2) = c(a^2 + b^2)$ ; les racines sont $a$ et $c$.
    Exemple : $2, \sqrt{20}, 10$, et $124x - x^3 = 240$ ;
  \item $n = 4, 5, 6$ : $B$ est la somme des puissances $(n-1)$-ièmes de $n$
    proportionnelles, $Z$ le produit d'un extrême par la somme de celles des
    autres, et les racines sont les deux extrêmes. Exemples :
    $2, \sqrt[3]{40}, \sqrt[3]{200}, 10$ et $1248x - x^4 = 2480$ ;
    $2, \sqrt[4]{80}, \sqrt{20}, \sqrt[4]{2000}, 10$ et
    $12496x - x^5 = 24960$ ; $2, \sqrt[5]{160}, \ldots, 10$ et
    $124992x - x^6 = 249920$.
\end{itemize}
Tout est juste, une fois corrigé le dernier second membre.\footnote{La page écrit
« $124{,}960$ », souligné, et la marge porte « $249920$ », qui est la bonne valeur :
$2\,(160 + 800 + 4000 + 20000 + 100000) = 249\,920$. Le second terme de la suite,
surchargé (« 360 » ou « 160 »), est récrit en marge « $LQC$ 160 », qui est juste.
Que ces équations aient exactement ces deux racines positives suit de la
syncrise : $Bx - x^n - Z$ s'annule aux deux extrêmes, et n'a pas d'autre racine
positive, par la règle des signes (voir la note sur le porisme de la vue 23).}

Pour les ambiguës « sous le degré le plus élevé », $Bx^{n-1} - x^n = Z$, $B$ est la
somme de $n$ proportionnelles, $Z$ le produit d'un extrême par la puissance
$(n-1)$-ième de la somme des autres, et les racines sont la somme des $n - 1$
premiers et la somme des $n - 1$ derniers. Sur la suite $1, 2, 4, \ldots$ :
$7x^2 - x^3 = 36$ (racines $3$ et $6$), $15x^3 - x^4 = 2744$ ($7$ et $14$),
$31x^4 - x^5 = 810\,000$ ($15$ et $30$), $63x^5 - x^6 = 916\,132\,832$ ($31$ et
$62$). Tout est juste.\footnote{Pour $n = 3$ : si $x = a + b$, alors
$x^2(B - x) = c\,(a + b)^2 = a\,(b + c)^2$, puisque $b + c = r(a + b)$ et $c = ar^2$.
Les « 7Q » et « 31QQ » de la vue 60 sont écrits « 72 » et « 3122 », le Q ayant la
forme d'un 2.}

Les « degrés intermédiaires » (vues 60 à 63) se ramènent aux précédents par une
racine plane ou solide.\footnote{Pour la troisième des propositions qui suivent, la
page écrit « $B$ plano compositum » là où la dimension demande un plan-plan. Pour
la dernière, elle donne « fit $1N$ 25 vel 100 », valeurs de $x^2$, et la marge
corrige en $5$ et $10$, valeurs de $x$. La page dit aussi que, prises comme
équations en $x^2$ ou en $x^3$, $17N - 1Q = 16$ a les racines $1$ et $16$,
$513N - 1Q = 512$ les racines $1$ et $512$, et $21N - 1C = 20$ les racines $1$ et
$4$.}
\begin{itemize}
  \item $Bx^2 - x^4 = Z$ : $B$ est la somme de deux carrés $u^2 + v^2$, $Z$ leur
    produit, et $x^2$ vaut l'un ou l'autre ; exemple $17x^2 - x^4 = 16$, racines $1$
    et $4$ ;
  \item $Bx^3 - x^6 = Z$ : $B = u^3 + v^3$, $Z = u^3v^3$ ; exemple
    $513x^3 - x^6 = 512$, racines $1$ et $8$ ;
  \item $Bx^2 - x^6 = Z$ : $B$ est la somme des carrés de trois plans
    proportionnels, $Z$ le produit d'un extrême par la somme des carrés des deux
    autres, et $x^2$ vaut un extrême ; exemple $21x^2 - x^6 = 20$, avec les plans
    $1, 2, 4$, racines $1$ et $2$ ;
  \item $Bx^4 - x^6 = Z$ : $B$ est la somme de trois plans proportionnels, $Z$ le
    produit d'un extrême par le carré de la somme des deux autres, et $x^2$ vaut
    la somme des deux premiers ou des deux derniers ; exemple, avec les plans
    $5, 20, 80$ : $105x^4 - x^6 = 50\,000$, racines $5$ et $10$.
\end{itemize}
La vue 62 montre comment trouver trois plans proportionnels
dont la moyenne, ajoutée au premier ou au dernier, fasse un carré : ce sont
\[ \frac{B^4}{B^2 + G^2}, \quad \frac{B^2G^2}{B^2 + G^2}, \quad
   \frac{G^4}{B^2 + G^2} , \]
car la moyenne plus le premier vaut $B^2$, plus le dernier $G^2$. Avec $B = 1$,
$G = 2$ on a $\tfrac15, \tfrac45, \tfrac{16}{5}$, et, multipliés par $25$, les
plans $5, 20, 80$ de l'exemple.\footnote{Aux dénominateurs, la page porte le double
trait, corrigé en « $+$ » dans la marge. Le troisième plan est écrit
« $\frac{16}{15}$ », le premier chiffre du dénominateur étant peut-être une rature ;
le calcul donne $\frac{16}{5}$. La marge établit la formule par
$(G - E)(D - E) = E^2$, d'où $E = GD/(G + D)$.}

Les deux dernières propositions du chapitre (« De reliquis », vues 62 et 63)
sont justes dans leurs exemples. Pour $Bx^2 - x^5 = Z$ et la suite $1, 2, 4$ :
$279x^2 - x^5 = 2268$, racines $3$ et $6$, sommes des deux premiers et des deux
derniers termes. Le coefficient vaut $B = u^3 + v^3 + a\,v^2$, où $u = a + b$,
$v = b + c$, et le terme connu $Z = u^2v^2(a + b + c)$.\footnote{La phrase qui
définit $B$ est lue avec des surcharges et ne donne pas, telle qu'elle est lue,
la valeur $279 = 27 + 216 + 36$ ; la note marginale qui définit $Z$ est raturée
et « très incomplète » selon la transcription. Les deux expressions données ici
sont de cette lecture ; elles se tirent de la syncrise,
$B = (u^5 - v^5)/(u^2 - v^2)$, et donnent bien $279$ et $9 \cdot 36 \cdot 7 = 2268$.}
Pour $Bx^3 - x^5 = Z$ et la suite $1, 2, 4, 8$ : $217x^3 - x^5 = 57\,624$, racines
$7$ et $14$ ; ici $B = u^2 + v^2 - c\,(a + b + c)$, et $Z = (B - u^2)\,v^3 =
(B - v^2)\,u^3$, avec $u = a + b + c$, $v = b + c + d$, ce que la page dit, une fois
intervertis par la marge les mots « primis » et « postremis ».\footnote{Vérifié :
$49 + 196 - 4 \cdot 7 = 217$ ; $(217 - 49) \cdot 343 = (217 - 196) \cdot 2744 =
57\,624$.}

\subsection*{Chapitre XIX : les contradictoires}

Si $x^2 + Bx = Z$ et $y^2 - By = Z$, les deux racines sont deux côtés dont la
différence est $B$ et le produit $Z$ ($x^2 + x = 2$ donne $1$, $y^2 - y = 2$
donne $2$). Si $x^4 + Bx = Z$ et $y^4 - By = Z$, il y a quatre proportionnelles
$a, b, c, d$ dont les cubes pris alternativement diffèrent de $B$,
$B = d^3 - c^3 + b^3 - a^3$, avec $Z = a\,(b^3 - c^3 + d^3)$, et $x = a$, $y = d$ :
sur $1, \sqrt[3]{2}, \sqrt[3]{4}, 2$, $x^4 + 5x = 6$ donne $1$ et $y^4 - 5y = 6$
donne $2$. De même au sixième degré avec six proportionnelles
($x^6 \pm 21x = 22$, racines $1$ et $2$). Si $x^4 + Bx^3 = Z$ et $y^4 - By^3 = Z$,
$B$ est la différence alternée des quatre termes, $x$ celle des trois premiers,
$y$ celle des trois derniers : sur $1, 2, 4, 8$, $x^4 \pm 5x^3 = 216$ donne $3$
et $6$. De même au sixième degré sur $1, 2, 4, 8, 16, 32$ : $x^6 \pm 21x^5 =
5\,153\,632$, racines $11$ et $22$.\footnote{La page écrit deux fois
« $4{,}618{,}812$ », souligné et surligné, et la marge donne « $5153632$ », qui est
juste : $11^6 + 21 \cdot 11^5 = 22^6 - 21 \cdot 22^5 = 5\,153\,632$. Le numéro du
chapitre est écrit « XVIII XIX ». La marge de ces pages, d'une écriture plus
serrée dont les transcriptions ne décident pas si elle est de l'écrivain,
démontre les règles sur la suite $A, B, B^2/A, B^3/A^2$ et vérifie que
$Z/x^3 = B + x$ et $Z/y^3 = y - B$ ; ces vérifications sont justes.}

\subsection*{Chapitres XX et XXI : les inverses}

Si $Bx - x^3 = Z$ et $y^3 - By = Z$, il y a trois proportionnelles dont les
carrés pris alternativement font $B = a^2 - b^2 + c^2$, avec $Z = a\,(c^2 - b^2)$,
$x = a$ et $y = c$ (sur $1, \sqrt{2}, 2$ : $3x - x^3 = 2$ donne $1$, $y^3 - 3y = 2$
donne $2$) ; de même au cinquième degré ($11x - x^5 = 10$ donne $1$,
$y^5 - 11y = 10$ donne $2$). Si $Bx^2 + x^3 = Z$ et $By^2 - y^3 = Z$, sur $1, 2, 4$ :
$3x^2 + x^3 = 4$ donne $1$ et $3y^2 - y^3 = 4$ donne $2$ ; au cinquième degré,
$11x^4 \pm x^5 = 10\,000$ donne $5$ et $10$. Avec six proportionnelles
$1, 2, \ldots, 32$ : $231x^3 \pm x^5 = 96\,040$ (racines $7$ et $14$) et
$297x^2 \pm x^5 = 2916$ (racines $3$ et $6$), où $231 = (1 + 32)(8 - 1)$ et
$297 = (1 + 32)(4 - 1)^2$.\footnote{Corrections de l'écrivain, toutes justes :
« $B$ cubus » corrigé en « $E$ » (vue 65) ; « $1QQ$ » corrigé en « $C$ », pour
$y^5 - 11y = 10$ (vue 66) ; au troisième théorème, un « $=$ » barré d'un trait
vertical, et « $+$ » en marge ; au cinquième, « maior » corrigé en « minor », et
« $231Q$ » corrigé deux fois en « $C$ » ; au sixième, « tertiam » corrigé en
« secundam ». Restent deux écarts de dimension : au cinquième théorème, $Z$ est
dit fait « du cube de la différence $\ldots$ dans $B$ planum plus le carré de la même
différence », là où l'exemple demande la puissance cinquième
($7^3 \cdot 231 + 7^5 = 96\,040$) ; au sixième, $Z$ est dit « planoplano » là où
l'énoncé a « planosolido ». Au quatrième, le signe manque devant « $E$
quadratocubo » dans la seconde équation.} Le chapitre XXI donne deux
constitutions de plus : si $Bx - x^3 = Z$ et $y^3 - By = Z$, sur trois
proportionnelles, $B = a^2 + b^2 + c^2$, $Z = (a + c)\,b^2$, $x$ vaut un extrême et
$y$ la somme des extrêmes ($21x - x^3 = 20$ donne $1$ ou $4$, $y^3 - 21y = 20$
donne $5$) ; et si $Bx^2 - x^3 = Z$ et $By^2 + y^3 = Z$, $B = a + b + c$,
$Z = a(b + c)^2$, $x = a + b$ ou $b + c$, $y = b$ ($7x^2 - x^3 = 36$ donne $3$ ou
$6$, $7y^2 + y^3 = 36$ donne $2$). C'est la fin du premier traité.\footnote{Dans la
dernière proposition, la première des deux constructions de $Z$ (« composita
autem e duabus primis \emph{adducta} compositæ e duabus reliquis, dum ducitur in
mediæ quadratum ») ne donne pas $36$ telle qu'elle est lue ; la seconde est
juste. Le mot en italique est une lecture douteuse. Le premier traité n'a pas de
conclusion : la vue 68 s'arrête sur cet exemple, la moitié inférieure blanche, et
le second traité commence en tête du feuillet suivant.}

\section*{Le « De Emendatione » : cinq préparations, et l'expurgation}

\pagerange{69}{79}
Le second traité commence en tête d'un nouveau feuillet, sous son titre, et dit
d'emblée ce qu'il fait :

\begin{quote}
« Agnita æquationum constitutione analysta ad præparandum eas quæ sunt alioqui
\emph{mechanicismo} respuant aut demum ægre subeunt toto se confert [\ldots]
fere autem quæ prosunt geometricæ ad εὐμηχανίαν prosunt et arithmeticæ ut
etiam e contra. [\ldots] nunc autem circa numerosam analysim magis esse intentum
nostri est instituti. »\footnote{« La constitution des équations une fois
reconnue, l'analyste se consacre tout entier à préparer celles qui, autrement,
refusent la résolution ou ne s'y prêtent qu'avec peine [\ldots] ; en général, ce
qui sert à la commodité de la résolution géométrique sert aussi à
l'arithmétique, et inversement. [\ldots] Mais notre dessein est ici de nous
appliquer surtout à l'analyse numérique. » Le mot en italique, et le mot grec
« εὐμηχάνους » un peu plus loin, sont des lectures douteuses de la
transcription ; « δυσμηχανία » (dusmêchania), dans le titre, est l'incommodité d'une
équation à se laisser résoudre, « εὐμηχανία » (eumêchania) sa commodité.}
\end{quote}

Les préparations « solennelles » sont au nombre de cinq, numérotées sur le
feuillet : l'expurgation par onces, la transmutation πρῶτον ἔσχατον
(prôton eschaton, « premier-dernier »),
l'anastrophe, l'isomérie et la paraplérose climactique. S'y ajouteront un
climactisme « symétrique » (chapitre V), l'hypostase double (chapitre VII) et
une transmutation « canonique » (chapitre VIII).

\subsection*{L'expurgation par onces}

C'est le remède contre la multiplicité des termes. Une équation dont figure le
terme de degré $n - 1$ vient, par translation, d'une équation qui ne l'a pas ;
pour le faire disparaître, on translate la racine en sens contraire d'une
« once » du coefficient : la moitié pour le carré (« diahemysis »), le tiers pour
le cube (« diatritemorion »), le quart pour le carré-carré, le cinquième, le
sixième, « et ainsi à l'infini ». Selon le signe du terme, on pose
$x = y - B/n$, $x = y + B/n$ ou $x = B/n - y$.\footnote{La page nomme ces trois cas
« tres casuum differentias ». Le nom de la première opération est lu
« Diahemysis » avec doute, et celui de la quatrième, « diatetartemorion », est
écrit dans la marge à la place de « et sic in infinitum ordine », biffé. Deux
mots de la même racine, vue 70 (« a \ldots, de \ldots »), ne sont pas lus.} En
termes d'aujourd'hui, c'est la réduction d'une équation à la forme
\emph{réduite}, sans second terme, par la substitution $x = y - a_{n-1}/n$.

Les deux exemples sont justes : $x^2 + 2Bx = Z$ devient, avec $y = x + B$,
$y^2 = Z + B^2$ ; et $x^3 + 3Bx^2 + Dx = Z$ devient, avec
$y = x + B$,\footnote{La page écrit d'abord « $+$ $B$ quadrato » au troisième terme ;
la marge corrige en « $- B$ quadrato ter », qui est juste. Le signe devant « $D$
plano in $B$ » est noirci, lu « $\pm$ » sans assurance ; le calcul donne $+$.}
\[ y^3 + (D - 3B^2)\,y = Z + DB - 2B^3 . \]
Ainsi,
conclut le chapitre, toutes les quadratiques se ramènent aux pures, les cubiques
aux cubiques affectées sous le seul côté, les quartiques à celles qui ne sont
affectées que sous le côté et le carré, et ainsi de suite.

\subsection*{Les « formules » de réduction}

Suit une phrase sur les anciens, qui s'interrompt au bas de la vue 72 :
« Credebant autem antiqui quoniam hac reductione æquationes quadraticas omnino
expurgabant et facile explicabant » ; la vue 73 reprend sur une ligne en partie
effacée, et la transcription ne sait pas si c'est la suite.\footnote{« Or les
anciens croyaient, parce que par cette réduction ils purgeaient entièrement les
équations quadratiques et les résolvaient facilement\ldots ». La vue 73 poursuit :
ils ont cherché obstinément à purger de même les degrés supérieurs, sans
chercher d'autre méthode, « et bonas horas mathematicas [\ldots] dispendio
absumpserunt » — ils y ont perdu de bonnes heures de mathématiques ; la méthode
n'est pas « catholica », générale. Plusieurs mots de ce passage sont lus avec
doute (« vel duobus », « accedere », « mechanicorum », « exquisierint »,
« erraverunt », « collocabant », « hæredium »).} Viennent des formules, chacune
avec son « conséquent » et des exemples :

\begin{itemize}
  \item $x^2 + 2Bx = Z$ : $x = \sqrt{Z + B^2} - B$ ; exemple $x^2 + 2x = 20$,
    $x = \sqrt{21} - 1$ ;
  \item $x^2 - 2Bx = Z$ : $x = \sqrt{Z + B^2} + B$ ; exemple
    $x = \sqrt{21} + 1$ ;
  \item $2Dx - x^2 = Z$ : $x = D \mp \sqrt{D^2 - Z}$ ; exemple $10x - x^2 = 20$,
    $x = 5 \mp \sqrt{5}$.
\end{itemize}

Et pour les cubiques affectées sous le carré, ramenées à des cubiques
affectées sous le côté :
\begin{itemize}
  \item $x^3 + 3Bx^2 = Z$, $y = x + B$ : $y^3 - 3B^2y = Z - 2B^3$ ; exemple
    $x^3 + 6x^2 = 1600$, $x = 10$, et $y^3 - 12y = 1584$, $y = 12$ ;
  \item $x^3 - 3Bx^2 = Z$, $y = x - B$ : $y^3 - 3B^2y = Z + 2B^3$ ; exemple
    $x^3 - 6x^2 = 400$, $x = 10$, et $y^3 - 12y = 416$, $y = 8$ ;
  \item $3Bx^2 - x^3 = Z$, $y = x - B$ ou $y = B - x$ :
    $3B^2y - y^3 = Z - 2B^3$ ou $2B^3 - Z$ ; exemples $21x^2 - x^3 = 972$,
    $x = 9$ ou $18$, et $147y - y^3 = 286$, $y = 2$ ou $11$ ; puis
    $9x^2 - x^3 = 28$, $x = 2$, et $27y - y^3 = 26$, $y = 1$.
\end{itemize}
Tout est juste.\footnote{« $1684$ » est corrigé en marge en « $1584$ ». Le premier
second membre se lit « 600 » ou « 1600 », le 1 mal séparé ; c'est $1600$. La
marge de la vue 74 ajoute « fit $1N$ 2 », la valeur de $B$, que le texte omet. Le
feuillet note qu'en arithmétique on peut marquer d'un signe les chiffres de la
racine altérée (« σημεῖον (sêmeion) aliquod superimponitur »). L'équation
$9x^2 - x^3 = 28$ a une seconde racine positive, $(7 + \sqrt{105})/2 \approx
8{,}62$, que la page ne donne pas.}

Puis sept formules pour les cubiques « affectées sous le carré et le côté »,
$x^3 \pm 3Bx^2 \pm Dx = Z$ et leurs inverses, toutes ramenées par $y = x \pm B$
ou $y = B - x$ à une cubique sans terme en $y^2$ ; ainsi, pour la première,
$x^3 + 3Bx^2 + Dx = Z$ et $y = x + B$,
\[ y^3 + (D - 3B^2)\,y = Z + DB - 2B^3 , \]
et de même pour les autres, les signes de la page étant souvent le double trait,
puisque l'ordre des grandeurs comparées dépend des données. Chacune porte deux à quatre
exemples ; tous ont été refaits et sont justes, une fois reçues les corrections
que l'écrivain a lui-même portées en marge.\footnote{Corrections marginales, toutes
justes : « $30N$ » en « $330N$ » (vue 75) ; « $76$ » en « $80$ » (vue 76) ; « $20$ »
en « $120$ » (vue 77) ; deux fois « $300Q$ » en « $30Q$ » (vues 77 et 78) ;
« $96$ » en « $36$ » (vue 78) ; « $330X$ » en « $330N$ » ; les signes laissés en
blanc ou doublés, rétablis en « $+$ » (vues 76, 77). Une faute n'est pas
corrigée : à la vue 78, « $30Q - 1C$ æquatur 56 », qui doit être $30x - x^3 = 56$
(racine $4$). Deux coefficients surchargés (« 24 », « 60 ») sont lus d'après le
calcul par la transcription.}

\subsection*{Les onces triangulaires et pyramidales}

La fin du chapitre (vues 78 et 79) étend l'expurgation : on peut aussi faire
disparaître, par une translation, un terme de degré plus bas, en prenant pour
once l'inverse d'un nombre triangulaire ou pyramidal. Pour $x^n$ affecté sous
$x^{n-2}$ par $B^2$, on prend le carré de la translation égal à $B^2$ divisé par
le nombre triangulaire $\binom{n}{2}$ (tiers pour le cube, sixième pour le
carré-carré, dixième pour le carré-cube) ; pour $x^n$ affecté sous $x^{n-3}$ par
$B^3$, le cube de la translation égal à $B^3$ divisé par le nombre pyramidal
$\binom{n}{3}$ (quart pour le carré-carré, dixième pour le carré-cube). C'est
juste : dans $(y + h)^n$, le coefficient de $y^{n-k}$ est $\binom{n}{k}h^k$. Le
traité ajoute, et c'est exact, qu'à la place du terme supprimé apparaissent les
autres.\footnote{Deux précisions que le traité ne donne pas. (1) Pour le terme
en $x^{n-2}$, il faut que ce terme soit retranché : $x^3 - B^2x$ se purge par
$h^2 = B^2/3$, mais $x^3 + B^2x$ demanderait $h^2 = -B^2/3$, qui n'est pas une
grandeur ; pour le terme en $x^{n-3}$, le signe est indifférent, puisqu'on
prend une racine cubique. (2) Faire disparaître un terme autre que le second en
créant les autres n'avance pas la résolution ; c'est pour en supprimer plusieurs
à la fois que Tschirnhaus (1683) prendra une substitution de degré plus élevé.
La marge corrige « $B$ cubo » en « $B$ quadrato » pour le carré-cube affecté sous le
cube, et « æqualitate » en « æqualitas ». La suite « nempe sunt numeri
triangulares 3. 6. 10. 15 \&c. » est en tête du feuillet suivant.}

\section*{La transmutation πρῶτον ἔσχατον (prôton eschaton)}

\pagerange{79}{82}
Le chapitre II est le remède contre la négation : quand le terme connu est
« nié » de la puissance, c'est-à-dire quand les termes affectés l'emportent, on
prend pour nouvelle inconnue le quotient du terme connu par la racine,
$y = Z/x$. Le nom vient de la proportion qui en résulte, où « le premier devient
le dernier ». Pour $x^3 - Bx = Z$, on obtient
\[ y^3 + By^2 = Z^2 , \]
où l'affection négative est devenue positive ; et, en effet,
$x : \sqrt[3]{Z} = \sqrt[3]{Z^2} : (x^2 - B)$, avec $y = x^2 - B$. Exemple :
$x^3 - 96x = 40$ ; $y = 40/x$ donne $y^3 + 96y^2 = 1600$, de racine $4$, d'où
$x = 10$.\footnote{La page écrit « $1C \,-\, 46Q$ æquabitur $160$ », souligné, et la
marge corrige en « $+ 96$ » et « $1600$ » : $4^3 + 96 \cdot 4^2 = 1600$. La page écrit
aussi « explicando » et « vocatur », corrigés en marge en « applicando » et
« renovatur », et « æquationum » en « adfectionum ».} La même transformation
fait disparaître certaines irrationnelles : $x^3 - 10x = \sqrt{48}$ devient
$y^3 + 10y^2 = 48$, de racine $2$, et $x = \sqrt{48}/2 = \sqrt{12}$.

Pour la quartique $x^4 - Bx^3 - Dx^2 = Z$, $y = Z/x$ donne
\[ y^4 + DZ\,y^2 + BZ^2\,y = Z^3 , \]
avec l'exemple $x^4 - 3x^3 - 8x^2 = 50$ : $y^4 + 400y^2 + 7500y = 125\,000$, de
racine $10$, d'où $x = 5$ ; et $x^4 - 8x = \sqrt[3]{80}$ donne $y^4 + 8y^3 = 80$,
de racine $2$, d'où $x = \sqrt[3]{10}$.\footnote{La page écrit « $12500$ »,
surligné, et la marge « $125000$ », qui est juste. La troisième ligne du calcul
en espèces porte deux fois « $Z$ plano plano » là où l'on attend « $Z$ plano plano
plano plano » ; le produit est le même.}

Le chapitre résume la transformation en quatre préceptes, justes : le terme
connu, inchangé pour les quadratiques, s'élève au carré pour les cubiques, au
cube pour les quartiques, et ainsi de suite ; les affections négatives
deviennent positives et inversement ; le coefficient du terme en $x^k$ se
multiplie par $Z^{k-1}$ ; et la racine cherchée est le terme connu divisé par la
nouvelle. Il remarque enfin qu'on pourrait diviser par une autre grandeur que
le terme connu, mais que celui-ci évite les fractions, et qu'en géométrie on
préfère une grandeur homogène au carré de la racine.\footnote{Les préceptes sont à
la vue 82, sur un petit feuillet sans numéro ; qu'ils soient ceux qu'annoncent
les quatre dernières lignes de la vue 81 est une inférence de la transcription.
Deux mots en sont lus avec doute (« Inuariato », « pronunciato »). L'exemple
final, $x^3 - Bx^2 = BZ^2$ avec $x = BZ/y$, donne bien $y^3 + B^2y = B^2Z$.}

\section*{L'anastrophe}

\pagerange{83}{92}
Le chapitre III traite les équations « inversement niées » de degré impair,
$Bx - x^3 = Z$, qui ont deux racines positives et que l'analyse numérique
résout mal. L'anastrophe les ramène à une équation plus basse par
l'intermédiaire de leur corrélative : si $y$ est la racine de $y^3 - By = Z$
(qui n'en a qu'une positive), on ajoute $y^3$ aux deux membres de
$x^3 = Bx - Z$ :
\[ x^3 + y^3 = Bx + By , \]
et l'on divise par $x + y$ :
\[ x^2 - xy + y^2 = B , \qquad\text{soit}\qquad yx - x^2 = y^2 - B . \]
Connaissant $y = D$, les deux racines de l'équation proposée sont celles de cette
quadratique. Exemple : $39x - x^3 = 70$ ; la corrélative $y^3 - 39y = 70$ a la
racine $7$ ; d'où $7x - x^2 = 10$, et $x = 2$ ou $5$.\footnote{La page écrit
d'abord « 3 vel 6 », biffé, et la marge corrige en « 2. vel. 5. », qui est juste.
C'est, en termes d'aujourd'hui, la division du polynôme $x^3 - Bx + Z$ par le
facteur $x + D$ que fournit sa racine négative $-D$ : la racine positive de la
corrélative est l'opposée de la racine négative de la proposée. Le traité, qui
ne connaît que les racines positives, obtient ce facteur par l'équation
corrélative.}

De même, pour $Bx^2 - x^3 = Z$, avec $y^3 + By^2 = Z$ de racine $D$ :
$(B + D)x - x^2 = D^2 + BD$. Exemple : $7x^2 - x^3 = 36$ ; $y^3 + 7y^2 = 36$ donne
$2$ ; d'où $9x - x^2 = 18$, $x = 3$ ou $6$.\footnote{Au théorème II, le signe
devant « $B$ » est le double trait, corrigé en « $+$ » dans la marge ; au
problème II, « fit regularis » est corrigé en « fiat irregularis », comme
« fit » en « fiat » au problème I.} Le traité dit que la méthode vaut pour les
degrés impairs, et renvoie les degrés pairs à la paraplérose climactique.

Au cinquième degré, pour $Bx - x^5 = Z$, avec $y^5 - By = Z$ de racine $D$ :
\[ D^3x - D^2x^2 + Dx^3 - x^4 = D^4 - B ; \]
exemple : $11x - x^5 = 10$ ; $y^5 - 11y = 10$ donne $2$ ; d'où
$8x - 4x^2 + 2x^3 - x^4 = 5$, qui donne $1$. Et pour $Bx^4 - x^5 = Z$, avec
$y^5 + By^4 = Z$ :
\[ (BD^2 + D^3)\,x - (BD + D^2)\,x^2 + (B + D)\,x^3 - x^4 = BD^3 + D^4 ; \]
exemple : $11x^4 - x^5 = 10\,000$ ; $y^5 + 11y^4 = 10\,000$ donne $5$ ; d'où
$400x - 80x^2 + 16x^3 - x^4 = 2000$, qui donne $10$. Tout est juste. Au cinquième
degré, cependant, l'anastrophe n'abaisse l'équation qu'au quatrième, et
l'exemple le montre : $11x - x^5 = 10$ a une seconde racine positive,
$x \approx 1{,}4026$, qui est aussi racine de la quartique obtenue, et que la
page ne donne pas.\footnote{Vérifié : $x^5 - 11x + 10 = (x - 1)(x + 2)(x^3 - x^2 +
3x - 5)$, dont le dernier facteur a la racine réelle $1{,}4026\ldots$. Les
marges de ces pages corrigent plusieurs signes (« $+$ » pour le double trait) et
ajoutent un terme omis (« $+ 1D$ in $A$ cubum », vue 89), toutes corrections justes.}

La seconde moitié du chapitre prend « la voie contraire » : si l'on connaît une
racine d'une équation ambiguë, on obtient les autres par la même division. Quatre
théorèmes, justes.\footnote{Au dernier théorème, la page divise d'abord par
« $A + D$ », que la marge corrige en « $A - D$ », qui est juste ; la division
suppose $x \ne D$. Au deuxième, « minus » est corrigé en « plus ». Le double
trait « $=$ » de ces pages (« $A$ cubus $=$ $D$ cubo », « per $A = D$ ») est la
différence d'ordre non donné entre les deux racines ; la marge le corrige en
« $+$ » là où c'est une somme.}
\begin{itemize}
  \item $x^3 - Bx = Z$ et $By - y^3 = Z$ de racine $D$ : $x^2 - Dx = B - D^2$ ;
    exemple $x^3 - 8x = 7$ avec $D = 1$ : $x = \sqrt{29/4} + \tfrac12$ ;
  \item $Bx^2 + x^3 = Z$ et $By^2 - y^3 = Z$ : $x^2 + (B - D)x = BD - D^2$ ;
    exemple $9x^2 + x^3 = 8$ avec $D = 1$ : $x = \sqrt{24} - 4$ ;
  \item $Bx - x^3 = Z$, deux racines $x$ et $D$ : $x^2 + Dx = B - D^2$ ; exemple
    $8x - x^3 = 7$ avec $D = 1$ : $x = \sqrt{29/4} - \tfrac12$ ;
  \item $Bx^2 - x^3 = Z$, deux racines $x$ et $D$ : $x^2 + (D - B)x = BD - D^2$ ;
    exemple $9x^2 - x^3 = 8$ avec $D = 1$ : $x = 4 + \sqrt{24}$.
\end{itemize}

\section*{L'isomérie, et le climactisme symétrique}

\pagerange{92}{97}
Le chapitre IIII, l'\emph{isomérie}, chasse les fractions : on réduit les
fractions au même dénominateur $D$, on pose $y = Dx$, et l'on multiplie le
coefficient du terme en $x^{n-1}$ par $D$, celui du terme en $x^{n-2}$ par $D^2$,
et ainsi de suite ; la racine nouvelle est la racine ancienne multipliée par le
dénominateur. Ou, inversement, on divise. La justification est qu'égaux
multipliés ou divisés par égaux donnent des égaux. Les exemples sont
justes.\footnote{Corrections de la marge, justes : le rapport des racines « ut 8 ad 1 »
en « 48 ad 1 », et « ut $\frac{18}{2}$ ad 1 » en « 12 ad 1 » ; « $+ B$ solido in $A$ »
rétabli à la vue 93 ; « hic » et « illic » intervertis à la vue 95. Un passage de
la vue 94 est barré d'une croix, avec en marge « penitus expungenda », et une
note encadrée de la vue 95, « sic restituantur », le récrit ; la transcription
en donne la lecture « douteuse dans le détail ». Une note marginale numérotée
« 4 » (vue 93) est reprise par « 4. vide annotata pag. 48 » à la vue 100. Le mot
lu « sumptis » à la vue 95 est la lecture la plus douteuse de la page.}
\begin{itemize}
  \item $x^3 + \tfrac32 x = 225$ devient $y^3 + 6y = 1800$, $y = 12$, $x = 6$ ;
    $x^3 + \tfrac32 x = \tfrac{265}{2}$ devient $y^3 + 6y = 1060$, $y = 10$,
    $x = 5$ ;
  \item $x^3 + \tfrac32 x^2 = 270$ devient $y^3 + 3y^2 = 2160$, $y = 12$ ;
    $x^3 + \tfrac32 x^2 = \tfrac{325}{2}$ devient $y^3 + 3y^2 = 1300$, $y = 10$ ;
  \item $x^3 + \tfrac{11}{12}x = \tfrac{19}{4}$ devient, avec $y = 48x$,
    $y^3 + 2112y = 525\,312$, $y = 72$, $x = \tfrac32$ ; ou, avec le dénominateur
    commun $12$, $y^3 + 132y = 8208$, $y = 18$ ; et l'on passe de l'une à l'autre
    en divisant par $4$ ;
  \item $x^3 + 12x^2 + 8x = 2280$, $x = 10$, divisée par $2$, donne
    $y^3 + 6y^2 + 2y = 285$, $y = 5$ ; et $x^3 = 1728$ donne $y^3 = 8$.
\end{itemize}

Le chapitre V, le \emph{climactisme symétrique}, chasse les irrationnelles en
élevant les deux membres à la puissance qui les fait disparaître. Si
$x^3 - Bx = \sqrt{Z}$, on élève au carré :
\[ x^6 - 2Bx^4 + B^2x^2 = Z , \]
une cubique en $x^2$ ; exemple : $x^3 - 2x = \sqrt{1200}$ devient, en $u = x^2$,
$u^3 - 4u^2 + 4u = 1200$, de racine $12$, et $x = \sqrt{12}$. Le traité note qu'on
y arrivait aussi par πρῶτον ἔσχατον : $y^3 + 2y^2 = 1200$, $y = 10$,
$x = \sqrt{1200}/10 = \sqrt{12}$. Si $x^3 - \sqrt[3]{B}\,x = Z$, on divise par $x$,
on élève au cube et l'on multiplie par $x^3$ :
\[ x^9 - 3Zx^6 + (3Z^2 - B)\,x^3 = Z^3 ; \]
exemple : $x^3 - \sqrt[3]{18}\,x = 6$ devient, en $u = x^3$,
$u^3 - 18u^2 + 90u = 216$, de racine $12$, et $x = \sqrt[3]{12}$. Tout est
juste.\footnote{La page écrit « $B$ plano in $A$ quad. » au deuxième terme du carré,
corrigé en marge en « $+ B$ plano-plano in $Aq$ », et « $8Q$ », corrigé en « $18$ ».
Le premier paragraphe du chapitre est d'une écriture rapide dont plusieurs mots
sont douteux (« facit », « desideres », « veram », « libata », « nec facta »), et
la transcription le dit non grammatical sous la forme lue.}

\section*{La paraplérose climactique : les quartiques par une cubique}

\pagerange{97}{119}
Le chapitre VI est le cœur du second traité. Son titre dit ce qu'il fait :
« Quemadmodum æquationes quadratoquadraticæ deprimuntur ad quadraticas per
medium cubicarum a radice plana. Seu de climactica paraplerosi. »\footnote{« Comment
les équations du quatrième degré s'abaissent aux quadratiques par le moyen
d'équations cubiques dont la racine est plane. Ou de la paraplérose
climactique. » « Paraplérose » est le « remplissage » : le chapitre dit que le nom
vient du supplément qu'on ajoute. Le paragraphe d'ouverture (vue 97) dit aussi
que les cinq modes suffisent à préparer les équations pour la résolution
numérique, que les racines irrationnelles s'y obtiennent « vero proximæ », au
plus près du vrai, et qu'« exhiber les racines exactes est l'affaire du géomètre
plutôt que de l'arithméticien ». Au pied de la vue 98, « regularis » est corrigé
en « irregularis ».}

\subsection*{La méthode}

Soit $x^4 + Bx = Z$. On écrit $x^4 = Z - Bx$ et l'on ajoute aux deux membres
$x^2y^2 + \tfrac14 y^4$, où $y$ est une grandeur auxiliaire : le premier membre
devient le carré $\bigl(x^2 + \tfrac12 y^2\bigr)^2$. On veut que le second,
$y^2x^2 - Bx + Z + \tfrac14 y^4$, soit aussi un carré, et l'on choisit pour sa
racine
\[ \frac{B}{2y} - yx , \]
dont le carré est $y^2x^2 - Bx + B^2/(4y^2)$. « Ainsi, dit le feuillet, les
affections sous $A$ disparaîtront dans la comparaison, et l'on sera ramené à une
équation en $E$, où l'on voulait aller. » Il faut et il suffit que
$Z + \tfrac14 y^4 = B^2/(4y^2)$, c'est-à-dire
\[ y^6 + 4Z\,y^2 = B^2 , \]
une équation du troisième degré en $y^2$. Si $D$ en est la racine, les deux
carrés sont égaux, leurs racines aussi, et
\[ x^2 + Dx = \frac{B}{2D} - \frac{D^2}{2} . \]
C'est la méthode que l'on appelle aujourd'hui de Ferrari : ajouter à une
quartique un terme qui en fasse une égalité entre deux carrés, et déterminer le
paramètre par une équation du troisième degré, la \emph{résolvante}.\footnote{Ferrari
l'avait trouvée et Cardan l'a publiée dans l'\emph{Ars magna} (1545) ; le nom de
résolvante est moderne. Le traité ne cite personne. Descartes (1637) donnera une
autre réduction de la quartique au troisième degré.} Tout est juste pour la
racine positive que cherche le traité ; la quartique a en outre une racine
négative et deux racines imaginaires, dont le traité ne parle pas.

Pour $x^4 - Bx = Z$, la racine à comparer est $B/(2y) + yx$, la résolvante est la
même, et $x^2 - Dx = B/(2D) - D^2/2$. Pour $Bx - x^4 = Z$, la résolvante est
$y^6 - 4Zy^2 = B^2$, et $Dx - x^2 = D^2/2 - B/(2D)$, ce qui suppose
$B/(2D) < D^2/2$, comme le note une marge de la vue 104.\footnote{Pour ce dernier
cas, on ajoute $x^2y^2 + \tfrac14y^4$ à $x^4 = Bx - Z$ ; la racine à comparer est
$yx + B/(2y)$. La note de la vue 104 est numérotée « 4 », et c'est à elle que
renvoie, à la vue 100, « 4. vide annotata pag. 48 » : la vue 104 est le
feuillet 48 de la série à l'encre.}

\subsection*{Les trois théorèmes, et la constitution}

La vue 101 et le béquet « 47 bis » (vues 102 et 103) en tirent trois théorèmes et
les vérifient par la syncrise : l'équation $x^4 + Bx = Z$ est faite de deux côtés
$x$ et $x + D$, avec
\[ B = 2x(x + D)\,D + D^3 , \qquad Z = x(x + D)\,\bigl(x^2 + xD + D^2\bigr) , \]
ou, sur quatre proportionnelles $a, b, c, d$, avec $D = d - a$,
$B = d^3 - c^3 + b^3 - a^3$ et $Z = a\,(b^3 - c^3 + d^3)$ — la constitution des
contradictoires du premier traité (chapitre XIX). L'exemple est juste : sur
$2, \sqrt[3]{40}, \sqrt[3]{200}, 10$, l'équation $x^4 + 832x = 1680$ a pour
résolvante, en $u = y^2$, $u^3 + 6720u = 692\,224$, de racine $64$ ; donc $D = 8$,
et $x^2 + 8x = 52 - 32 = 20$ donne $x = 2$. De même, $x^4 - 832x = 1680$ donne
$x^2 - 8x = 20$, $x = 10$ ; et $1248x - x^4 = 2480$ a pour résolvante
$u^3 - 9920u = 1\,557\,504$, de racine $144$, d'où $D = 12$ et $12x - x^2 = 20$,
$x = 2$ ou $10$.\footnote{« $1557$ », surligné dans le texte, est donné en entier
dans la marge, « $1{,}557{,}504$ », le dernier chiffre surchargé : $1248^2 =
1\,557\,504$. Au théorème II, « $A$ solido-solido » est corrigé en « $B$ », et
« $-E$ quadrato » en « $D$ quadrato ½ ». Dans les constitutions, $D$ est la
différence des deux côtés aux théorèmes I et II, leur somme au théorème III, où
$B$ est la somme des cubes des quatre proportionnelles et $Z$ le produit d'un
extrême par la somme des cubes des trois autres.}

\subsection*{La quartique affectée sous le cube}

Le problème II traite $x^4 + 2Bx^3 = Z$. On forme le carré de
$x^2 + Bx - \tfrac12 y$, où $y$ est un plan :
\[ \bigl(x^2 + Bx - \tfrac12 y\bigr)^2 = x^4 + 2Bx^3 + (B^2 - y)\,x^2 - Byx
   + \tfrac14 y^2 . \]
Il faut que $Z + (B^2 - y)x^2 - Byx + \tfrac14 y^2$ soit un carré, de racine
$\sqrt{B^2 - y}\;x - By/\bigl(2\sqrt{B^2 - y}\bigr)$ ; d'où la résolvante
\[ y^3 + 4Z\,y = 4ZB^2 , \]
et, $D$ en étant la racine,
\[ x^2 + \Bigl(B + \sqrt{B^2 - D}\Bigr)x = \frac{D}{2} + \frac{DB}{2\sqrt{B^2 - D}} . \]
Pour $x^4 - 2Bx^3 = Z$, les signes de $B$ changent ; pour $2Bx^3 - x^4 = Z$, la
résolvante est $y^3 - 4Zy = 4ZB^2$ et la racine carrée porte sur $B^2 + y$. Les
exemples sont justes : sur $1, 2, 4, 8$, $x^4 + 5x^3 = 216$ a pour résolvante
$y^3 + 864y = 5400$, de racine $6$ ; comme $\sqrt{4B^2 - 4D} = \sqrt{25 - 24} = 1$,
on a $x^2 + 3x = 18$, $x = 3$ ; $x^4 - 5x^3 = 216$ donne $x^2 - 3x = 18$, $x = 6$ ;
$15x^3 - x^4 = 2744$ a pour résolvante $y^3 - 10976y = 617\,400$, de racine $126$,
d'où $\sqrt{225 + 504} = 27$ et $21x - x^2 = 98$, $x = 7$ ou $14$.\footnote{Le
feuillet écrit « $\frac{1}{2}$ » de trop au terme « $-E$ plano in $A$ quad. », et le
barre à la ligne suivante ; « quadratum » est corrigé en « cubus » dans la
résolvante, et le coefficient $4$ rétabli en marge. Il écrit la racine
« $LV\{B\,q - E\,p\}$ in $A\,q$ », la racine universelle portant sur tout le produit,
soit $\sqrt{(B^2 - y)\,x^2} = \sqrt{B^2 - y}\,x$ : la marge corrige « in $A$ » en
« in $A\,q$ » dans ce sens. Aux exemples, « $- 5C$ » est corrigé en « $+$ », « $540$ »
en « $5400$ », « $725$ » en « $729$ ». Ici « $B$ bis » vaut $5$, puis $15$ : la
différence alternée, puis la somme, des quatre proportionnelles.}

\subsection*{La quartique complète, par deux formules}

Le problème III traite $x^4 \pm 2Gx^2 \pm Bx = Z$, la quartique privée seulement
de son terme en $x^3$ — la forme à laquelle l'expurgation ramène toute
quartique. La « première formule » forme le carré de
$x^2 + G + \tfrac12 y^2$ et compare au carré de $B/(2y) - yx$ ; la résolvante est
\[ y^6 + 4G\,y^4 + (4Z + 4G^2)\,y^2 = B^2 , \]
et $x^2 + Dx = B/(2D) - D^2/2 - G$ (avec les signes changés selon les cas). La
« seconde formule » (« Aliter », vues 111 à 113), pour $x^4 + Gx^2 + Bx = Z$,
ajoute $Fx^2 + \tfrac14F^2$, où $F$ est un plan, et obtient la résolvante
\[ F^3 - G\,F^2 + 4Z\,F = B^2 + 4ZG , \]
puis $x^2 + \sqrt{F - G}\;x = B/\bigl(2\sqrt{F - G}\bigr) - F/2$. Chaque formule
donne sept théorèmes, un par combinaison de signes, chacun avec un exemple ; on
les a tous refaits. Par exemple, $x^4 + 6x^2 + 880x = 1800$ a pour résolvante,
selon la première formule ($G = 3$), $u^3 + 12u^2 + 7236u = 774\,400$, de racine
$64$, et, selon la seconde ($G = 6$), $F^3 - 6F^2 + 7200F = 817\,600$, de racine
$70$, avec $70 - 6 = 64$ ; dans les deux cas $x^2 + 8x = 20$, $x = 2$. De même
$44x^2 + 720x - x^4 = 1600$ donne $u = 144$ (première formule) ou $F = 100$,
$100 + 44 = 144$ (seconde), puis $12x - x^2 = 20$, $x = 2$ ou $10$. Tous les
exemples sont justes, une fois reçues les corrections de la marge ; une faute
reste dans l'énoncé du quatrième théorème de la seconde formule.\footnote{Ce
théorème, pour $x^4 - Gx^2 - Bx = Z$, écrit la résolvante avec
« $+ Z$ plano-plano $4$ in $G$ plano » au second membre ; il faut
$F^3 + GF^2 + 4ZF = B^2 - 4ZG$, comme au théorème III, où la marge a corrigé le
signe, et comme l'exige l'exemple, qui est le même
($F^3 + 4F^2 + 6400F = 614\,400 = 640\,000 - 25\,600$). Au même exemple, « $64N$ »
est écrit pour « $6400N$ ». Les corrections marginales de ces pages, toutes
justes, sont nombreuses : signes rétablis (« $+$ », « $\pm$ », « $\mp$ » à la place du
double trait), coefficients $4$ ajoutés, « Aq » pour « A » sous la racine
universelle, et des nombres : « $20$ » en « $10$ », « $160$ » en « $1600$ » (deux
fois), « $6410$ » en « $6416$ », « $800$ » en « $880$ », « $6$ » en « $60$ »,
« $36$ » en « $30$ ». Au théorème V de la seconde formule, la ligne
« $1C = 44Q = 6400N$ æquatur $800{,}000$ », aux signes doublés, se lit
$F^3 + 44F^2 - 6400F = 800\,000$. Les « équations après métamorphose » que la
marge donne pour chaque cas des vues 107 à 113 sont les résolvantes ; elles sont
justes.}

\subsection*{La feuille volante « 46 bis »}

Une feuille volante posée sur la vue 99, numérotée « 46 bis » d'une autre main,
récapitule en deux tableaux tous les cas : pour chaque forme de quartique, la
« racine » à former et la grandeur à lui comparer. On les a vérifiés ligne à
ligne ; tous sont justes, sauf un signe.\footnote{À la quatrième ligne du tableau
de gauche (série I à VII), pour $x^4 - 2Gx^2 - Bx = Z$, la racine doit être
$B/(2y) + yx$, comme à la ligne II ; la feuille écrit « $-$ E in A ». Dans le
tableau de droite, trois écarts d'écriture : « Eq4 » pour « Ep4 » (ligne VI),
« Eq ½ » pour « Ep ½ » (ligne V), et un dénominateur sans « Lv » (ligne VII). Une
bande de papier sous le pied de la vue 100, de la même main, porte deux
constitutions en espèces que la transcription ne rattache à rien de sûr.}

\subsection*{Deux notes marginales}

La marge de la vue 100 porte, en tête d'une addition encadrée, les mots
« transcripta ex exemplari » — lecture probable, non certaine — puis un calcul
abrégé du problème ; à quoi ces mots s'appliquent, la marge ne le dit pas. Et, au
problème IIII de la première formule (vue 114), au-dessus de la résolvante, des
mots soulignés puis biffés se lisent, avec beaucoup de doute, « deerant hæc in
exemplari ». Si ces lectures sont bonnes, l'écrivain copie un modèle et note où
son modèle est défaillant. Cette lecture ne va pas plus loin.\footnote{« Copiée
sur l'exemplaire » ; « ceci manquait dans l'exemplaire ». La transcription de la
vue 114 écrit : « S'ils se lisent bien ainsi, ils disent que le passage manquait
dans un modèle : la note serait celle d'un copiste. Cette transcription ne
tranche pas. » Le chapitre se clôt à la vue 119 par une phrase de lecture
difficile sur les métathèses et l'expurgation, puis « Hæc itaque sunto satis
superque ».}

\section*{L'hypostase double : la formule de Cardan}

\pagerange{119}{122}
Le chapitre VII revient à la cubique, avec la septième manière de changer
d'inconnue. Pour $x^3 + 3Bx = 2Z$, on pose
\[ y^2 + xy = B , \qquad\text{c'est-à-dire}\qquad x = \frac{B - y^2}{y} , \]
et l'équation devient
\[ y^6 + 2Z\,y^3 = B^3 , \]
une équation du second degré en $y^3$, « de racine solide ». D'où le
conséquent : si $D^3 = \sqrt{B^3 + Z^2} - Z$, alors $x = (B - D^2)/D$. Avec
$y^2 - xy = B$, on trouve de même $y^6 - 2Zy^3 = B^3$, $G^3 = \sqrt{B^3 + Z^2} + Z$
et $x = (G^2 - B)/G$. Et, les deux ensemble,
\[ x = G - D = \sqrt[3]{\sqrt{B^3 + Z^2} + Z} - \sqrt[3]{\sqrt{B^3 + Z^2} - Z} . \]
C'est la formule de Cardan, et la substitution est celle qu'on appelle
aujourd'hui la substitution de Viète.\footnote{L'énoncé coïncide avec la formule
publiée par Cardan en 1545 pour $x^3 + px = q$, avec $p = 3B$ et $q = 2Z$ ; le
chapitre VI du premier traité en contenait la constitution. Le nom de
substitution de Viète est moderne.} Les exemples sont justes :
$x^3 + 81x = 702$ ; $\sqrt{27^3 + 351^2} = \sqrt{142\,884} = 378$ ;
$378 - 351 = 27 = 3^3$, d'où $x = (27 - 9)/3 = 6$ ; et $378 + 351 = 729 = 9^3$,
d'où $x = (81 - 27)/9 = 6$ ; enfin $x^3 + 6x = 2$ donne
$x = \sqrt[3]{4} - \sqrt[3]{2}$.\footnote{« $700$ », surligné, est corrigé en
« $702$ » dans la marge ; « $B$ plano ter » est complété en « ter in $A$ » ; la
marge justifie le conséquent en complétant le carré de $y^3 + Z$.}

Le problème II traite $x^3 - 3Bx = 2Z$ sous la condition que le cube du tiers
du coefficient soit moindre que le quart du carré du terme connu, $B^3 < Z^2$ —
c'est la bonne condition. Avec $xy - y^2 = B$, $x = (B + y^2)/y$, l'équation
devient $2Zy^3 - y^6 = B^3$, et
\[ x = \sqrt[3]{Z + \sqrt{Z^2 - B^3}} + \sqrt[3]{Z - \sqrt{Z^2 - B^3}} ; \]
exemple : $x^3 - 81x = 756$, $\sqrt{378^2 - 27^3} = 351$, $x = 9 + 3 = 12$. C'est
juste. Le feuillet remarque que la condition découle de la forme même de
l'équation réduite ; quand elle n'est pas remplie, on est dans le cas
irréductible, que le premier traité résolvait par les sections angulaires
(vues 16 à 18).\footnote{Une note marginale d'encre plus noire ajoute que
« les mêmes choses suivent » pour $3Bx - x^3 = 2Z$. Ce n'est pas le cas pour une
racine positive : sous la condition $Z^2 > B^3$, cette équation n'a pas de racine
positive, puisque $3Bx - x^3$ ne dépasse pas $2B^{3/2}$ ; la formule y donne la
valeur absolue de sa racine négative. La marge écrit d'ailleurs « $Zs.$ » là où
l'énoncé porte « $Z$ solido bis ». Une autre note marginale (vue 122) justifie la
règle : $2Z$ est la somme de deux cubes dont le produit est $B^3$.}

\section*{La transmutation canonique ; les racines entières}

\pagerange{122}{125}
Le chapitre VIII apprend à transformer une équation pour que son coefficient,
ou son terme connu, soit une grandeur prescrite. On pose $x = By/X$ pour
remplacer le coefficient $B$ par $X$ :
\[ x^3 + Bx^2 = Z \;\Longrightarrow\; y^3 + Xy^2 = \frac{X^3 Z}{B^3} , \]
et de même pour les autres formes ; pour prescrire le terme connu, on pose
$x = Zy/D$. Les exemples sont justes : $x^3 + 20x^2 = 96\,000$ devient
$y^3 + y^2 = 12$, $y = 2$, $x = 40$ ; $x^3 - 144x = 10\,368$ devient $y^3 - y = 6$,
$y = 2$, $x = 24$ ; $x^3 + 860x = 12^3$ devient $y^3 + 215y = 6^3$, $y = 1$,
$x = 2$.\footnote{La marge corrige deux dénominateurs, « $X$ » en « $X\,q.$ » et
« $Z$ solido » en « $X\,c$ in $Z\,s.$ sur $B\,c$ », tous deux justement. « Concedetur »
est lu par le sens (« Concreditor » ?).}

Le chapitre le justifie par une remarque de calcul : quand le coefficient vaut
l'unité, on voit tout de suite si la racine est entière. Pour $x^3 + x = 10$, le
cube le plus proche est $8$, et $8 + 2 = 10$ : la racine est $2$. Pour
$x^3 - x = 24$, on trouve $3$. Pour $x^3 + x = 9$ ou $x^3 - x = 25$, le calcul
tombe à côté, « donc la racine est irrationnelle ». La conclusion est juste,
mais elle suppose un fait que le chapitre ne dit pas : une racine rationnelle
d'une équation à coefficients entiers dont le premier coefficient vaut $1$ est
un entier.\footnote{C'est le théorème des racines rationnelles, appliqué au
cas unitaire : si $p/q$ irréductible est racine, $q$ divise le premier
coefficient. Il est postérieur au traité. Ici, $x^3 + x$ est croissante et prend
les valeurs $2$ et $10$ en $1$ et $2$ : aucun entier ne donne $9$ ; de même
$x^3 - x$ vaut $24$ et $60$ en $3$ et $4$.}

\section*{Réductions « anomales »}

\pagerange{125}{135}
Les chapitres IX et X rassemblent des équations qu'on abaisse d'un degré parce
qu'on leur connaît un facteur. Le traité n'en fait pas une règle : « de
irregularibus autem non statuitur præcepta quoniam anomalia non magis finita
est quam artificis in indagando vis et solertia ».\footnote{« On ne donne pas de
préceptes pour les irrégulières, parce que l'anomalie n'a pas plus de bornes que
la force et l'adresse de l'artisan à chercher. »} En termes d'aujourd'hui, chaque
proposition divise le polynôme par $x \pm B$, où $\mp B$ est une racine, ce que
l'on appelle le théorème du facteur ; et chaque fois que le facteur divisé est
$x - B$ avec $B$ positif, la division fait perdre la racine $x = B$, que le
traité ne donne pas toujours.

\subsection*{Chapitre IX}

\begin{itemize}
  \item $x^3 - 2B^2x = B^3$ donne $x^2 - Bx = B^2$ (on ajoute $B^3$ et l'on divise
    par $x + B$) ; exemple $x^3 - 18x = 27$, $x^2 - 3x = 9$. La racine est
    $x = \tfrac{1 + \sqrt5}{2}\,B$ : $B$ et $x$ sont dans la \emph{moyenne et extrême
    raison}.
  \item $2B^2x - x^3 = B^3$ donne $x^2 + Bx = B^2$ (on divise par $x - B$), dont la
    racine $\tfrac{\sqrt5 - 1}{2}\,B$ est le plus grand segment de $B$ coupé en
    moyenne et extrême raison ; mais $x = B$ est aussi racine, et la division la
    perd.\footnote{La page donne pour exemple « $18N - 1C$ æquatur $27$, igitur
    $1Q - 3N$ æquabitur $9$ » ; la règle qui précède donne $1Q + 3N$, et c'est
    juste : $x^3 - 18x + 27 = (x - 3)(x^2 + 3x - 9)$, de racines positives $3$ et
    $\tfrac{-3 + \sqrt{45}}{2} \approx 1{,}854$. La marge corrige « $B$ quadratum »
    en « bis ».}
  \item $x^3 - 3B^2x = 2B^3$ a la racine $x = 2B$ ; exemples $x^3 - 12x = 16$,
    $x = 4$, et $x^3 - 3x = 2$, $x = 2$. Et $3B^2x - x^3 = 2B^3$ a la racine double
    $x = B$ : $12x - x^3 = 16$ donne $2$, $6x - x^3 = \sqrt{32}$ donne
    $\sqrt2$.\footnote{La page énonce d'abord la proposition III sous la forme
    « $B$ quadrat. in $A$ ter minus $A$ cubo æquetur $B$ cubo bis, $B$ dupla est
    ipsius $A$ » ; la marge, d'encre plus noire, la récrit sous la forme juste
    ($A^3 - 3B^2A = 2B^3$, $A$ double de $B$) et ajoute une proposition IIII
    ($3B^2A - A^3 = 2B^3$, $A = B$). Le second exemple de la page, « $1C - 6N$ æquat.
    $L\,32$ fit $1N$ $L\,2$ », est faux pour la proposition III (la racine de
    $x^3 - 6x = \sqrt{32}$ est $2\sqrt2 = \sqrt8$) et juste pour la proposition
    IIII, d'où il paraît venir. En effet $x^3 - 3B^2x - 2B^3 = (x - 2B)(x + B)^2$ et
    $3B^2x - x^3 - 2B^3 = -(x - B)^2(x + 2B)$.}
  \item $x^3 - Bx^2 + Dx = BD$ a la racine $B$ ($x^3 - 4x^2 + 5x = 20$, $x = 4$) ;
    $x^3 + Bx^2 - D^2x = BD^2$ a la racine $D$ ($x^3 + 5x^2 - 4x = 20$, $x = 2$) ;
    $Bx^2 + D^2x - x^3 = D^2B$ a les racines $B$ et $D$ ($6x^2 + 4x - x^3 = 24$,
    $x = 6$ ou $2$). Toutes sont justes : ce sont $(x - B)(x^2 + D) = 0$,
    $(x + B)(x^2 - D^2) = 0$ et $(x - B)(x^2 - D^2) = 0$.\footnote{« $B$ cubus » est
    corrigé en « $D$ » à la proposition VI.}
  \item $Dx^2 + BDx - x^3 = B^3$ donne $(D + B)x - x^2 = B^2$ ; exemple
    $10x^2 + 20x - x^3 = 8$, $12x - x^2 = 4$, $x = 6 \pm \sqrt{32}$.\footnote{La page
    écrit « $1Q + 20N - 1C$ », corrigé en marge en « $10$ » ; le signe devant
    « $B$ in $A$ » est d'abord le double trait, puis un « $+$ » net.}
  \item $x^3 - Dx^2 + DBx = B^3$ donne $(D - B)x - x^2 = B^2$ ; exemple
    $8x - x^2 = 4$, $x = 4 \pm \sqrt{12}$. Ici encore la division par $x - B$ fait
    perdre la racine $x = B$.\footnote{L'exemple est écrit « $1C + 10Q + 20N$
    æquetur $8$ » ; la proposition demande $x^3 - 10x^2 + 20x = 8$, qui a les trois
    racines $2$, $4 - \sqrt{12}$ et $4 + \sqrt{12}$.}
  \item $x^3 - 3Bx = \sqrt{2B^3}$ a la racine $x = \bigl(\sqrt{3B} + \sqrt B\bigr)/\sqrt2$ ;
    exemple $x^3 - 6x = 4$, $x = \sqrt3 + 1$. Et $3Bx - x^3 = \sqrt{2B^3}$ a les
    racines $\bigl(\sqrt{3B} - \sqrt B\bigr)/\sqrt2$ et $\sqrt{2B}$ ; exemple
    $6x - x^3 = 4$, $x = \sqrt3 - 1$ ou $2$. Une marge vérifie le premier cube par
    un développement juste.\footnote{À l'exemple de la proposition X, la page
    écrit « igitur $L\,3 - 1$ fit $1N$ » ; la marge porte « $\pm$ » ; la racine positive
    de $x^3 - 6x = 4$ est $\sqrt3 + 1$, et $\sqrt3 - 1$ est celle de la
    proposition XI. Le second membre « $L\,B$ plano-plano-plani $\frac{\cdot}{1}$ » est
    corrigé en marge en « $2$ ». En effet $x^3 - 6x - 4 = (x + 2)(x^2 - 2x - 2)$.}
\end{itemize}

\subsection*{Chapitre X}

Le chapitre X, « Similium reductionum Continuatio », multiplie une quadratique
par $x \pm B$ ou $B - x$ :
\begin{itemize}
  \item $x^3 + 3Bx^2 + Dx = 2B^3 - DB$ donne $x^2 + 2Bx = 2B^2 - D$ ; exemple
    $x^3 + 30x^2 + 44x = 1560$, $x^2 + 20x = 156$, $x = 6$ ;
  \item $x^3 + 3Bx^2 - Dx = 2B^3 + DB$ donne $x^2 + 2Bx = 2B^2 + D$ ; exemple
    $x^3 + 30x^2 - 24x = 2240$, $x^2 + 20x = 224$, $x = 8$ ;
  \item $x^3 - 3Bx^2 + Dx = DB - 2B^3$ donne $2Bx - x^2 = D - 2B^2$, et la racine
    $x = B$ ; exemple $x^3 - 30x^2 + 236x = 360$, racines $2$, $10$ et $18$ ; et
    $x^3 - 30x^2 + 264x = 640$, racines $4$, $10$, $16$ ;
  \item $3Bx^2 + Dx - x^3 = 2B^3 + DB$ donne $x^2 - 2Bx = 2B^2 + D$, et la racine
    $B$ ; exemple $30x^2 + 24x - x^3 = 2240$, racines $28$ et $10$ ;
  \item $3Bx^2 - Dx - x^3 = 2B^3 - DB$ donne $x^2 - 2Bx = 2B^2 - D$, et la racine
    $B$ ; exemple $30x^2 - 156x - x^3 = 440$, racines $22$ et $10$.
\end{itemize}
Tout est juste, et le chapitre, cette fois, donne toujours la racine $B$ que la
multiplication a introduite.\footnote{Corrections de la marge, justes : « $756$ »
en « $156$ », « $22400$ » en « $2240$ », « $- D$ plano » rétabli, « æquabitur » en
« æquetur », plusieurs signes rétablis à la place du double trait. À la troisième
proposition, la condition lue « sit $B$ quadratum ter maius $D$ plano »,
$3B^2 > D$, est celle qui assure que la quadratique a des racines réelles ; pour
qu'elles soient toutes deux positives il faut en outre $D > 2B^2$, ce que
l'exemple vérifie ($200 < 236 < 300$). Le numéro du chapitre est écrit « XX »,
ramené à « X » dans la marge, comme les suivants.}

La proposition VI, « ad quadrato-quadraticam pertinens », est d'une autre
sorte : pour $x^4 - 2Xx^3 + 4X^3x = 2X^4$, la page, corrigée par la marge, trouve
\[ 2X^2x^2 - x^4 = \bigl(4 - 2\sqrt3\bigr)\,X^4 , \]
d'où $x^2 = X^2 \pm X^2\sqrt{\sqrt{12} - 3}$ ; exemple $X = 1$,
$x^4 - 2x^3 + 4x = 2$. Seul le signe moins donne une racine de l'équation
proposée : $x = \sqrt{1 - \sqrt{\sqrt{12} - 3}} \approx 0{,}5646$. Le signe plus
donne $1{,}2966$, valeur absolue de la racine négative, qui est racine de
$x^4 + 2x^3 - 4x = 2$ et non de l'équation proposée.\footnote{Vérifié : les
racines réelles de $x^4 - 2x^3 + 4x - 2$ sont $0{,}5646$ et $-1{,}2966$, et toutes
deux satisfont $2x^2 - x^4 = 4 - \sqrt{12}$. Le passage intermédiaire, où la page
élève au carré pour retrouver une équation « quadrato-cubique », a des
coefficients qui ne répondent pas à l'égalité précédente, et plusieurs mots n'en
sont pas lus ; il n'est pas restitué. La marge corrige « $X$ quad. in $A$ quad. »
en « $2$ », « $X$ quad.-quad. » en « $4$ », et « $L\,X$ quadrati » en « $3$ ».}

\section*{Quatre recueils ; les relations entre coefficients et racines}

\pagerange{135}{141}
Les chapitres XI à XIIII sont des recueils de « constitutions singulières »,
sans démonstration : « hæc singula suam habent ex Zetesi demonstrationem ».

\subsection*{La moyenne et extrême raison, « duplicata, triplicata\ldots »}

Le chapitre XI part de $x^2 + Bx = B^2$ : la longueur $2B$ est coupée en trois
segments proportionnels, $B$, $x$ et $B - x$, et l'on dit que $B$ est coupée en
moyenne et extrême raison ; $x = \tfrac{\sqrt5 - 1}{2}B$. Il généralise :
$x^3 + Bx^2 + B^2x = B^3$ coupe $2B$ en quatre segments continûment
proportionnels $B, x, x^2/B, x^3/B^2$ (« moyenne et extrême raison
doublée »), et ainsi de suite jusqu'à sept segments. C'est juste : la somme des
segments est $2B$ si et seulement si $x^n + Bx^{n-1} + \cdots + B^{n-1}x = B^n$.\footnote{Le
rapport $B/x$ est alors la racine réelle plus grande que $1$ de
$t^n = t^{n-1} + \cdots + t + 1$ : le nombre d'or pour $n = 2$, et, pour les degrés
suivants, les nombres qu'on appelle aujourd'hui « tribonacci »,
« tétranacci »\ldots{} Ces noms sont modernes. À la cinquième proposition, la page
écrit « $B$ in $A$ quad.-quad. » là où la suite demande $B^2x^4$. Le numéro du
chapitre est corrigé de « XXI » en « XI ».}

Le chapitre XII fait de même avec $B$ premier terme d'une suite de
proportionnelles et $Z$ la somme des autres : $x^2 + Bx = BZ$,
$x^3 + Bx^2 + B^2x = B^2Z$, \ldots, où $x$ est le second terme. Le chapitre XIII
prend les sommes alternées : $Bx - x^2 = BZ$, avec deux racines, le second terme
et la différence du premier et du second ; $x^3 - Bx^2 + B^2x = B^2Z$, etc. Les
exemples sont justes : sur $4, 6, 9$, $9x - x^2 = 18$ donne $6$ et $3$, et
$x^2 - 4x = 12$ donne $6$ ; sur $1, 2, 4, 8$, $x^3 - 8x^2 + 64x = 192$ donne $4$ ;
sur $1, 2, 4, 8, 16$, $4096x - 256x^2 + 16x^3 - x^4 = 20\,480$ donne $8$ ; sur
$1, 2, \ldots, 32$, $x^5 - 32x^4 + 1024x^3 - 32768x^2 + 1048576x = 11\,534\,336$ donne
$16$.\footnote{Au chapitre XIII, proposition II, le signe devant « $B$ in $A$
quad. » est le double trait, corrigé en « $-$ » dans la marge. Les marges des vues
137 et 138, d'encre plus noire, énoncent d'autres propositions de la même
série ; elles ne sont pas vérifiées ici, et l'une d'elles omet un terme
($+B\,x^3$ dans la somme des quatre derniers). Une autre note en tête de la
vue 138, dans un encadré, n'est pas lue assez pour être interprétée.}

\subsection*{Le chapitre XIIII : les équations à plusieurs racines}

Le dernier chapitre, « Collectio quarta », « laisse librement à l'examen
analytique » des théorèmes sur les équations « admirablement explicables par
plusieurs racines ». Ce sont les relations entre coefficients et racines, dans
leur forme générale :
\begin{itemize}
  \item $(B + D)\,x - x^2 = BD$ a pour racines $B$ et $D$ ;
  \item $x^3 - (B + D + G)\,x^2 + (BD + BG + DG)\,x = BDG$ a pour racines $B$, $D$
    et $G$ ;
  \item $(BDG + BDH + BGH + DGH)\,x - (BD + BG + BH + DG + DH + GH)\,x^2
    + (B + D + G + H)\,x^3 - x^4 = BDGH$ a pour racines $B$, $D$, $G$, $H$ ;
  \item et de même au cinquième degré, pour cinq racines $B$, $D$, $G$, $H$, $K$.
\end{itemize}
Les exemples sont justes : $3x - x^2 = 2$ (racines $1$, $2$) ;
$x^3 - 6x^2 + 11x = 6$ ($1, 2, 3$) ; $50x - 35x^2 + 10x^3 - x^4 = 24$
($1, 2, 3, 4$) ; $x^5 - 15x^4 + 85x^3 - 225x^2 + 274x = 120$ ($1, 2, 3, 4,
5$). C'est l'énoncé que l'on appelle aujourd'hui les formules de Viète : le
polynôme unitaire de racines $r_1, \ldots, r_n$ a pour coefficients, au signe près,
les fonctions symétriques élémentaires $\sigma_1 = \sum r_i$,
$\sigma_2 = \sum_{i<j} r_ir_j$, \ldots, $\sigma_n = r_1\cdots r_n$, de sorte que
\[ x^n - \sigma_1x^{n-1} + \sigma_2x^{n-2} - \cdots + (-1)^n\sigma_n = 0 . \]
Le traité l'énonce pour des racines positives données, et répartit les termes
entre les deux membres pour qu'ils soient tous positifs ; les signes alternés
que l'on écrit aujourd'hui y sont déjà.\footnote{L'énoncé coïncide exactement, pour
des racines positives, avec les relations de Viète. Leur forme avec des racines
quelconques, négatives et imaginaires comprises, est chez Girard (1629), et la
démonstration générale suppose le théorème fondamental de l'algèbre ; ce sont
des développements postérieurs au traité, si celui-ci est de Viète, comme le
disent ses titres. Corrections de la marge : « $- A$ qq. », omis à la proposition
III, et « $16$ » corrigé en « $10$ » dans son exemple ; à la vue 139, le titre
« prop. III » est biffé et la proposition récrite à la page suivante. La
proposition IIII s'achève en tête de la vue 141, dont le dernier membre porte le
double trait et un mot lu « æqu\ldots » seulement.}

Le traité s'achève là, sur une phrase qui dit ce qu'il pense de ce qu'il vient de
faire :

\begin{quote}
« Atq; hæc \emph{elegans} et perpulchra Speculationis Sylloge tractatui alioquin
effuso finem aliquem et coronida tandem Imponito. »\footnote{« Et que ce recueil
élégant et très beau de spéculation mette enfin un terme et un couronnement à un
traité par ailleurs trop étendu. » Le premier mot en italique est une lecture
douteuse : seules les lettres « -gans » sont nettes sous le cachet. Suit
« Finis », souligné. Un lecteur postérieur a écrit dans la marge « pag. 158
editionis elzevir. anni 1646, infolio. ».}
\end{quote}

\section*{Les tables des prostaphérèses de la Lune}

\pagerange{142}{147}
Les feuillets d'astronomie commencent à la vue 142, en tête du feuillet 67, sous
le titre « Ad Tabulas Prostaphæreseōn Lunarium ». La \emph{prostaphérèse} est
l'angle qu'il faut ajouter au mouvement moyen, ou en retrancher, pour avoir le
mouvement vrai — ce qu'on appelle aujourd'hui une \emph{équation}, au sens
astronomique. Les grandeurs dont elle dépend sont l'anomalie, la longitude
« double » de la Lune au Soleil, c'est-à-dire deux fois l'élongation moyenne,
qu'on note ici $2\eta$, et les distances des centres, exprimées en parties d'une
« altitude moyenne ».\footnote{Rien sur ces feuillets ne dit d'où viennent les
paramètres. Le petit tableau de la vue 145, « Canonici numeri in Sole », les
donne pour le Soleil et pour la Lune : $AB = 700\,000\,000$, $BC = 3\,000\,000$,
$EG = 27\,000\,000$, $GI = 2\,000\,000$ au Soleil ; $AB = 9\,100\,000$,
$BC = 100\,000$, $EG = 1\,000\,000$, $GI = 200\,000$ à la Lune. Le titre ne nomme que
le Soleil.}

Cinq propositions, chacune avec un exemple calculé pour $2\eta = 60^\circ$ :

\begin{itemize}
  \item \textbf{La première prostaphérèse.} Le centre $E$ de l'épicycle décrit,
    autour d'un point $B$ situé à la distance moyenne $AB$ de la Terre $A$, un
    petit cercle, l'« eccentrulus », de rayon $BC = AB/91$, et il y avance selon
    $2\eta$. L'angle $EAB$ est la première prostaphérèse, additive dans le premier
    demi-cercle de $2\eta$, soustractive dans le second :
    \[ \tan EAB = \frac{BE\sin 2\eta}{AB + BE\cos 2\eta} . \]
    Pour $2\eta = 60^\circ$ : $EF = 86\,602$, $FB = 50\,000$, $AF = 9\,150\,000$ en
    parties où $AB = 9\,100\,000$, et l'angle vaut $0^\circ\,32\tfrac{15}{29}'$. C'est
    juste ; le feuillet calcule par la tangente, « ut $AF$ ad $EF$ ita $100\,000$ ad
    $946$ ».
  \item \textbf{L'altitude moyenne corrigée} est $AE$, l'hypoténuse du triangle
    $AFE$. La page donne $9\,150\,131$ ; le calcul donne
    $\sqrt{9\,150\,000^2 + 86\,602^2} \approx 9\,150\,410$.\footnote{La page appelle cette
    grandeur « $EF$ », là où le raisonnement porte sur l'hypoténuse $AE$. L'écart de
    $279$ parties sur neuf millions ne change pas la cinquième proposition au
    degré de précision où elle calcule.}
  \item \textbf{La seconde prostaphérèse.} Un second petit cercle, de centre $G$ et
    de rayon $GI$, porté à la distance $EG$, est parcouru selon $2\eta$ à partir de
    son point le plus bas ; l'angle $GEK$ est la seconde prostaphérèse. Pour
    $2\eta = 60^\circ$, avec $EG = 1\,000\,000$ et $GK = 200\,000$ : $KL = 173\,205$,
    $EL = 900\,000$, et l'angle vaut $10^\circ\,53\tfrac35'$, ce qui est juste
    ($\tan = 0{,}19245$).\footnote{Une ligne interlinéaire dit « talium 1. qualium $GE$
    10. », alors que les nombres de l'exemple et du tableau de la vue 145 sont dans
    le rapport de $1$ à $5$. La page, écrite en recopiant la proposition I puis
    corrigée, garde des traces des nombres de celle-ci, biffés. Elle écrit
    « $AG$ » là où la figure a « $EG$ », et « $EXK$ » pour le triangle $KLE$.}
  \item \textbf{Le demi-diamètre corrigé de l'épicycle} est
    $EK = \sqrt{900\,000^2 + 173\,205^2} = 916\,515$. Juste.
  \item \textbf{La prostaphérèse maximale} a pour sinus $EK/AE$, soit
    $916\,515/9\,150\,131 = 0{,}10016$, et vaut $5^\circ\,44\tfrac{26}{29}'$. Juste :
    un épicycle de rayon $r$ vu à la distance $R$ donne une équation maximale
    $\arcsin(r/R)$.\footnote{La page écrit « ut $AE$ $AK$ ita Sinus Totus ad Sinum
    maximæ prostaphæreseos », $AK$ pour $EK$. Les nombres biffés d'une première
    rédaction ($953\,939$, $10\,425$) sont eux aussi cohérents entre eux. L'angle
    lui-même n'est pas écrit sur cette page ; il l'est dans la table de la vue 147.}
\end{itemize}

La vue 147 porte, tracés à la plume, les cadres des deux tables, celle de la
Lune de $0$ à $60$ degrés de $2\eta$ par pas de cinq, et celle du Soleil ; seule la
ligne de $60^\circ$ de la table lunaire est remplie, avec les cinq valeurs
ci-dessus.

En termes d'aujourd'hui, ce modèle fait deux choses. La première prostaphérèse
est une inégalité en longitude à peu près proportionnelle à $\sin 2\eta$, nulle
aux syzygies et aux quadratures, maximale aux octants, d'amplitude
$\arcsin(1/91) \approx 0^\circ 38'$ : c'est la forme de l'inégalité que l'on
appelle la \emph{variation}. La seconde fait dépendre de $2\eta$ le rayon effectif
de l'épicycle, $EK$, qui va de $800\,000$ à $1\,200\,000$ parties, et donc
l'équation maximale de l'épicycle : c'est le dispositif par lequel Copernic
rendait la seconde inégalité de la Lune, que l'on appelle aujourd'hui
l'\emph{évection}.\footnote{La variation (amplitude moderne $39'30''$, en
$\sin 2D$, où $D$ est l'élongation) a été découverte par Tycho Brahe. L'évection
(amplitude moderne $1^\circ 16'$, en $\sin(2D - M)$, où $M$ est l'anomalie) est la
seconde inégalité de Ptolémée ; Copernic la rendait par un petit épicycle qui
fait varier le rayon du premier avec $2D$, comme ici. Le mot d'évection est
postérieur à Ptolémée comme à Copernic, et il n'est pas sur ces feuillets. Ces
noms sont donnés parce que les dépendances coïncident ; les amplitudes du
feuillet ne sont pas comparées ici à celles d'aucun auteur ancien.}

\section*{L'hypothèse lunaire de l'« astroscope danois »}

\pagerange{148}{156}
Les vues 148 à 152 contiennent un brouillon très raturé (vue 148), puis sa mise au
net (vues 150 à 152), d'un texte qui commence par un jugement : l'hypothèse
lunaire de Tycho Brahe est mal construite, et l'écrivain la « restitue ». Le
brouillon nomme Tycho ; la mise au net biffe le nom et écrit à la place « ab
Astroscopo Dano », « l'astroscope danois » :

\begin{quote}
« Hypothesis Lunæ ab Astroscopo Dano proposita [\ldots] eum Geometram nequaquam
[\ldots] arguit. An igitur suis fidendum observatis Viderint \emph{αὐτόπται} (autoptai). Ego
suam Hypothesim ita restituo. »\footnote{« L'hypothèse de la Lune proposée par
l'astroscope danois montre qu'il n'est nullement géomètre. Faut-il donc se fier à
ses observations ? Que ceux qui ont vu de leurs yeux en jugent. Moi, je restitue
son hypothèse ainsi. » Le mot grec, écrit dans la marge à la place de
« Astroscopi » biffé, se lit « αὐλοσδι » (aulosdi) à la lettre ; « αὐτόπται » (autoptai) est la lecture
que la transcription propose par le sens. Le brouillon de la vue 148 portait
« infelicitatem artificis sive in Geometricis sive Opticis et physicis manifesto
arguit », biffé.}
\end{quote}

\subsection*{Le modèle}

Huit articles le décrivent (vue 150) :
\begin{itemize}
  \item la Lune se meut sur un « ellipsidium » (une petite ellipse), dont le centre
    parcourt une ellipse, dont le centre est porté par l'« homocentre », le cercle
    de rayon moyen autour de la Terre ;
  \item le centre de l'ellipse oscille (« libratur ») sur une ligne droite, vers
    l'occident puis vers l'orient, selon la longitude double $2\eta$ ;
  \item l'ellipse est décrite selon l'anomalie $M$ comptée de l'apogée « in
    antecedentia » ; l'ellipsidium selon l'« anomalie mixte » $2\eta - M$, comptée
    du périgée « in consequentia » ;
  \item dans l'une et l'autre, le demi-axe transversal (vers la Terre) est le tiers
    du demi-axe droit (perpendiculaire) ;
  \item en parties où l'homocentre vaut $276$ : ellipse $10$ et $30$, ellipsidium
    $2$ et $6$, cercle de libration $3$.\footnote{L'article 6, qui énonce le rapport
    triple, est entouré de mots et de fragments de vers biffés ou ajoutés —
    « tergemina », « triformis », « tergeminamque Hecaten, tria virginis ora
    Dianæ », que la transcription reconnaît comme un vers de l'\emph{Énéide} (IV,
    511) ; le jeu porte, semble-t-il, sur « tripla » et la triple Hécate.}
\end{itemize}

L'exemple (vues 151 et 152) est calculé pour $M = 45^\circ\,37\tfrac12'$ et
$2\eta = 230^\circ\,23\tfrac25'$, d'où l'anomalie mixte
$184^\circ\,45\tfrac{9}{10}'$. En parties où $AB = 27\,600\,000$, la page trouve
$AI = 27\,600\,000 + 699\,360 + 199\,308 = 28\,498\,668$ pour la composante le long
de la ligne moyenne, et, perpendiculairement, $IH = 714\,800 - 16\,616 = 698\,184$,
triplé en $IK = 2\,094\,552$, auquel s'ajoute la libration $KP = 231\,114$, soit
$IP = 2\,325\,666$ ; l'angle $IAP$, différence entre le lieu moyen et le lieu
apparent, a pour tangente $0{,}08161$ et vaut $4^\circ\,40'$. Tout est juste :
$699\,360$ et $714\,800$ sont $10^6\cos M$ et $10^6\sin M$ (à la dernière unité
près), $199\,308$ et $16\,616$ sont $2 \cdot 10^5\cos$ et $2 \cdot 10^5\sin$ de
$4^\circ\,45\tfrac{9}{10}'$, et $231\,114$ vaut $3 \cdot 10^5 \sin 50^\circ\,23\tfrac25'$.\footnote{La
page reporte « $2\,325\,665$ » dans la proportion finale, pour $2\,325\,666$. Le
brouillon de la vue 148 contient les mêmes nombres, avec des colonnes de calcul et
une légende des lettres ; le transcripteur ne les a pas tous vérifiés, et ils ne
le sont pas ici.}

En termes d'aujourd'hui, et au premier ordre dans les petits rapports
$30/276$, $6/276$, $3/276$, le lieu de la Lune s'écarte donc du lieu moyen de
trois termes périodiques : un terme en $\sin M$ d'amplitude $30/276$ radian,
soit $6^\circ 14'$ ; un terme en $\sin(2\eta - M)$ d'amplitude $6/276$ radian,
soit $1^\circ 15'$ ; un terme en $\sin 2\eta$ d'amplitude $3/276$ radian, soit
$37'$. Ce sont les arguments de l'\emph{équation du centre}, de l'\emph{évection} et
de la \emph{variation}, les trois plus grandes inégalités de la longitude de la
Lune.\footnote{Les amplitudes modernes sont $6^\circ 17'$, $1^\circ 16'$ et $39'30''$.
Le rapprochement est de cette lecture, au premier ordre seulement : le modèle du
feuillet est géométrique, et ses signes dépendent des sens de comptage qu'il
fixe. Il ne dit rien de la date du texte, et cette lecture n'en tire aucune.}

\subsection*{Les problèmes pour construire les tables}

Sous le titre « Ad condendum Canonas Problematia », l'écrivain calcule les
inégalités de la nouvelle hypothèse dans des cas simples :
\begin{itemize}
  \item la première prostaphérèse due au cercle de libration, de rayon
    $1/92$ de la distance moyenne (vue 149) : pour $2\eta = 60^\circ$, le rapport est
    $941$ pour $100\,000$ ; la page garde la valeur $0^\circ\,32\tfrac{15}{29}'$ de la
    vue 142, calculée pour $946$ ; avec $941$, on trouve $0^\circ\,32\tfrac13'$ ;
  \item la prostaphérèse quand la Lune est à $45^\circ$ du Soleil (vue 152), par une
    figure où l'on reporte les longueurs $2$ et $3$ de l'ellipsidium et du
    cercle ;
  \item quand l'anomalie vaut $90^\circ$ ou $270^\circ$ (vue 153), avec un tableau
    des cas intitulé « Factio » et des différences en lettres grecques, que la
    transcription ne reconstitue pas ;
  \item à la pleine et à la nouvelle Lune (vue 155) : en parties où
    $AB : BF = 34\tfrac12 : 1$, pour une anomalie de $30^\circ$, $AE = 3\,536\,602$,
    $EG = 3\,EF = 150\,000$, et l'angle $EAG$ vaut $2^\circ\,25\tfrac34'$ ;
  \item à la demi-Lune (vues 155 et 156) : $AB : BF = 23 : 1$, $AE = 2\,386\,602$,
    et l'angle vaut $3^\circ\,35\tfrac45'$.
\end{itemize}
Les nombres des deux derniers problèmes sont justes.\footnote{Aux deux, la page
appelle l'angle trouvé « complementum AEG » ; c'est l'angle $EAG$ que l'énoncé
demande, et c'est lui qui vaut $2^\circ\,25\tfrac34'$ et $3^\circ\,35\tfrac45'$. Le
rapport $34\tfrac12$ à la pleine Lune et $23$ à la demi-Lune fait varier, selon la
phase, la distance au centre de l'ellipse : c'est encore le dispositif de
l'évection. La vue 154 porte deux figures soignées, « Sol. » et « Luna. ». La vue
156 porte encore l'énoncé biffé d'un autre problème, un tableau des
prostaphérèses aux phases écrit en travers de la marge, en lettres grecques, un
hexagramme inscrit dans un cercle, et une note effacée où la transcription lit,
avec doute, le mot français « monsieur » suivi d'un nom dont elle ne voit que
l'initiale R.}

\section*{L'« harmonique ptolémaïque », Copernic, et les planètes}

\pagerange{157}{163}
La vue 157 ouvre un chapitre, « Harmonici Ptolemaici Constructio CAPUT XII »,
dont on n'a que le préambule. Il commence par Copernic :

\begin{quote}
« Copernicus sibi suarumque hypotheseōn si qua est præstantiæ detrahit adeò eas
mala construit Geometriâ. Itaque cum in quinque Planetis angulum exquirit quem
vocat Prostaphæreseos eccentri, Triangulum rectangulum ex quo is statim dabatur
proscindit in duo, quæ dicuntur obliquangula sex septemue Analogias pro una
exercens. »\footnote{« Copernic retire à ses hypothèses, et à lui-même, ce
qu'elles ont de mérite, tant il les construit par une mauvaise géométrie. Ainsi,
quand il cherche dans les cinq planètes l'angle qu'il appelle de la
prostaphérèse de l'excentrique, il coupe en deux triangles, qu'on dit
obliquangles, le triangle rectangle qui le donnait tout de suite, employant six
ou sept proportions pour une seule. » La suite dit que ce « solécisme » est répété
cinq fois, et qu'Aristarque lui-même le pardonnerait à peine ; puis « Quamobrem
ea emenda \emph{χρεώστης} nouæ est » (le mot grec, lu « χρεωστες », chreôstes, est
une lecture douteuse ; le mot χρεώστης, chreôstês, désigne le débiteur), et « je construirai d'abord l'harmonique de Ptolémée, afin que, de la
comparaison entre celui de Ptolémée et celui de Copernic, la doctrine des
mouvements soit mieux visible de toutes parts ».}
\end{quote}

L'objection est de méthode et elle est juste dans son principe : dans les modèles
à excentrique, l'angle cherché se lit directement dans un triangle rectangle
dont on connaît deux côtés, comme le fait la première proposition des vues 142
et 143, par une seule proportion (la tangente). La page ne montre pas le calcul
de Copernic qu'elle critique, et cette lecture ne le compare à aucun imprimé.

La vue 158 donne, sans titre, une légende et une construction pour la Lune —
homocentre, ellipse, ellipsidium, petits cercles « ad ellipsim » et « ad
ellipsidium » — où le rapport des deux demi-axes est double, et non triple comme
à la vue 150.\footnote{La page conclut « tota $IQ$ dupla ad totum $IH$ » à partir de
$HR = 2\,QS$ et $IR = 2\,IS$ ; ces prémisses donnent $IH = 2\,IQ$, et la conclusion,
telle qu'elle est lue, inverse le rapport. Que cette page continue le chapitre XII
de la vue 157 est une inférence de la transcription.} La vue 160 s'ouvre au
milieu d'une phrase qui définit « l'angle $KAG$, parallaxe de l'orbe », puis
énonce une « Propositio II » pour « les mêmes planètes » — Saturne, Jupiter, Mars
et Vénus, précise la vue 161 : connaissant les anomalies, la prostaphérèse
apparente, l'altitude moyenne et les « adlaterales » vers la Terre, trouver
l'« adlateralis » vers le Soleil. La construction commence : un eccentrulus de
rayon $\varphi$ (le demi-axe transversal), un point $E$ pris selon l'anomalie
$\sigma$, la perpendiculaire $EF$ prolongée en $G$ avec $FG = 2\,EF$, et $G$ centre
de l'épicycle ; la page s'arrête là.\footnote{Au premier ordre, un centre
d'épicycle qui s'écarte de la ligne moyenne de $2\varphi\sin\sigma$ latéralement et
de $\varphi\cos\sigma$ en distance est ce que donnent, avec une excentricité
$\varphi$, l'équant de Ptolémée comme le mouvement képlérien. Le rapprochement est de
cette lecture ; la page ne dit rien de tel. La lettre de l'« adlateralis » est lue
$\varphi$ avec doute ; « emecata » (vue 161) n'est pas un mot connu de la
transcription.}

Les deux dernières vues (162 et 163), sous le titre grec « Δεδομένον » (Dedomenon, « donné »),
posent qu'à partir de l'altitude moyenne de Mercure, des « adlaterales » et des
anomalies à la Terre et au Soleil, on trouve la distance de Mercure au Soleil.
Suivent une « préparation » par proportions de sinus — l'altitude moyenne est
allongée (« protractio ») ou raccourcie (« contractio ») selon le quadrant de
l'anomalie double à la Terre, puis « corrigée ad Omœopathiam » — et quatre
« factæ », chacune une proportion « comme le sinus total au sinus de tel angle,
ainsi telle ligne à telle autre ». Le texte s'arrête sur « Ac denique Vt Altitudo
media Mercurij », au milieu d'une proportion ; un mot, « emendata », d'une encre
plus noire au pied de la dernière page, est peut-être une réclame. La suite n'est
pas dans le volume.\footnote{La page dit la protraction « in primo \& ultimo
quadrante », la contraction « in tertio \& quarto » : le dernier quadrant est
nommé deux fois ; on attend, pour la contraction, le deuxième et le troisième.
Les troisième et quatrième « factæ » nomment le même angle $IMO$ ; ainsi sur la
page. Les lettres $M$, $N$, $O$ n'ont pas été vues sur la figure de la vue 160 ; les
proportions ne sont donc pas vérifiées ici. « Omœopathiam » est ainsi écrit.}

\section*{Ce que ces feuillets ne disent pas}

\pagerange{1}{163}
\textbf{Ils ne disent pas qui les a écrits.} Deux titres, de la main du texte,
portent le nom de François Viète ; la couverture aussi. Les transcriptions
n'identifient aucune main, et cette lecture n'en identifie aucune. Une seconde
main, dans les marges du chapitre VI du premier traité, parle de ses propres
démonstrations sur les sections angulaires ; elle ne se nomme pas. Les notes
« transcripta ex exemplari » (vue 100) et « deerant hæc in exemplari » (vue 114),
et celle de la vue 20, pourraient faire de ces traités la copie d'un
modèle ; ce sont des lectures douteuses, et elles ne disent ni quel modèle, ni
qui copie. Les feuillets d'astronomie ne portent pas de nom ; ils nomment Tycho
Brahe, Copernic, Ptolémée et Aristarque, et ne sont attribués ici à
personne.\footnote{La transcription du lot 8 décrit l'écriture des feuillets
d'astronomie comme « la même cursive régulière » que celle du traité du lot 2 ;
c'est une description d'écriture, qu'elle ne pousse pas jusqu'à une attribution,
et cette lecture non plus.}

\textbf{Ils ne sont pas datés.} Le catalogue ne donne pas de datation ; cette
lecture n'en infère aucune, ni des titres, ni du nom de Tycho, ni des renvois
à l'édition de 1646, qui sont d'un lecteur. Chaque fois qu'un résultat est dit
« postérieur », c'est par rapport au traité si celui-ci est de Viète, comme
le disent ses titres ; les feuillets eux-mêmes peuvent être plus tardifs.

\textbf{Ils ont été lus sans imprimé.} Aucune édition imprimée des œuvres de
Viète n'a été ouverte : ni pour combler une lacune, ni pour trancher une leçon,
ni pour vérifier les renvois aux pages de 1646. Ce qui manque à la transcription
manque ici, et est dit à sa place.

\textbf{Ce qui reste illisible ou incertain.} Le théorème du feuillet B (vue 3),
dont l'énoncé ne se laisse pas reconstituer ; une proportion de la construction
du feuillet A (vue 2) ; les longues marges de la seconde main aux vues 16 à 18,
données « sous réserve » ; plusieurs mots grecs (vues 5, 32, 83, 150, 157) ; des
phrases entières de la syncrise (vues 47 et 54) et du climactisme symétrique
(vue 96) ; les calculs marginaux des vues 42, 55 et 138 ; la liaison entre les
vues 72 et 73 ; l'exemple « $1C + 9Q - 24N$ æq. $33$ » de la vue 32 ; le passage
de la proposition VI du chapitre X de l'« Emendatio » où l'équation est élevée au
carré (vue 134) ; les tableaux en lettres grecques des vues 153 et 156, et la note
française de la vue 156 ; les proportions des « factæ » de Mercure (vue 163), dont
la figure n'a pas été entièrement vue. Tout cela est signalé en note à sa place.

\end{document}
